\documentclass[11pt]{article}
\usepackage[a4paper,margin=25mm]{geometry}
\usepackage{amsmath,amssymb,amsthm}
\usepackage[british]{babel}
\usepackage[round,authoryear]{natbib}
\setlength{\bibhang}{1.5em}
\setlength{\bibsep}{2pt}
\renewcommand{\bibfont}{\small}
\usepackage{url}
\usepackage{hyperref}
\newtheorem{theorem}{Theorem}
\newtheorem{proposition}[theorem]{Proposition}
\newtheorem{lemma}[theorem]{Lemma}
\newtheorem{corollary}[theorem]{Corollary}
\newcommand{\E}{\mathbb E}
\newcommand{\Prb}{\mathbb P}
\title{Supporting certificates for guarded corrections and coherent pricing preparation}
\author{Amanjeet Singh}
\date{11 October 2026}
\begin{document}
\maketitle
\begin{abstract}
I construct and strengthen a matched quantum pricing rival which encodes a
nonnegative capped-payoff difference on a shared Gaussian draw. Guarded
arithmetic and explicit projection preserve the probability range when
distinct approximate tables fail to preserve endpoint order. I derive a
common-information stability certificate for positive Gaussian corrections,
separating fixed endpoint shifts from normal-dependent scaling remainders.
The same information is supplied to raw-payoff, numerical conditional-price
and shared-threshold implementations. All 416 tested certificates tighten,
while canonical gate caps remain unchanged. Complete reflection costs yield
an all-degree criterion for transferring preparation savings to parity
samplers. A matched protected Chebyshev comparison retains finite confidence,
discarded pilots and full fallbacks across the pricing alternatives.
The financial target remains the stipulated classical no-arbitrage model.
Established smoothing and reflection mechanisms are distinguished from the
checked graph-specific derivations and synthetic resource comparisons;
generic mathematical priority and physical quantum advantage are unestablished.
\end{abstract}
\smallskip\noindent\textit{Keywords:} quantum derivative pricing; common random numbers;
capped basket options; Gaussian discretisation; reversible arithmetic.\par
\smallskip\noindent\textit{JEL classification:} C63, C65, G13.\par

\section{Exact financial target}

Fix the outer factor coordinate $y$ in the classical risk-neutral model,
and let $m=m_G(y)$, $\Delta=\log T(y)\geq0$ and $\lambda>0$.
The ratio inequality follows from the positive basket weights and exact
weighted-centred drift and loading identities. Those premises and the
finite outer box, its grid transfer and analytic coefficient $c_0$ are
retained from the main paper. The geometric-conditioning and conditional
smoothing foundations are established financial methods
\citep{curran,bayer}; the retained main-paper source ledger records the
versions and sections previously consulted. Put $K=100$ and $H=1$, and define
\begin{equation}
 f(x)=\min\{1,\max\{100(e^x-1),0\}\},\qquad
 h_{m,\Delta}(n)=f(m+\Delta+\lambda n)-f(m+\lambda n).
\end{equation}
The monotonicity of $f$ gives $0\leq h_{m,\Delta}\leq1$ pointwise.
For $N\sim N(0,1)$,
\begin{equation}
 R_1(y)=\E h_{m,\Delta}(N)
       =P_\lambda(m+\Delta)-P_\lambda(m).
\end{equation}
Both terms use the same normal label. The financial model is unchanged;
only the coherent representation of its first correction changes.

\section{Exact narrow-domain payoff implementation}

\begin{lemma}[Clamp decomposition]\label{lem:raw-clamp}
Let $b=1/64$ and $x_+=\max(x,0)$. For every real $x$,
\begin{equation}
 f(x)=\min\{1,100\max\{\exp(\min\{b,x_+\})-1,0\}\}.
\end{equation}
\end{lemma}
\begin{proof}
Let $\ell=\log(101/100)$. Since $\log(1+u)<u$ for $u>0$,
$0<\ell<1/100<1/64$. If $x\leq0$, both expressions vanish.
If $0<x<\ell$, neither clamp changes $x$, so both expressions are
$100(e^x-1)$. If $\ell\leq x\leq b$, this value is at least one
and the final clamp returns one. If $x>b$, the exponential argument
is $b>\ell$, and the final clamp again returns one. Endpoints follow
from the same equalities. Thus no exponential outside $[0,b]$ is needed.
\end{proof}

At fixed precision $q=28$, each constant, addition, positive subtraction,
scaling and clamp uses the retained exact reversible integer emitter.
The nonnegative exponential on $[0,b]$ uses certified rounded ordinates
and signed-magnitude linear interpolation. The exact smooth interpolation
error is at most $e^b b^2/(8M^2)$ for $M=2^{b_e}$ cells, in addition
to certified nodal rounding and one interpolation floor unit.
The exact exponential output range is bounded by $[1,64/63]$:
for $0\leq x\leq1/64$, the power series gives
$e^x=\sum_{k\geq0}x^k/k!\leq\sum_{k\geq0}x^k=1/(1-x)\leq64/63$.
This avoids gratuitous fixed-point headroom.
No derivative is taken across a payoff kink: both kinks are implemented
by exact monotone clamps. The retained graph certificate propagates
all intermediate rounding errors and proves that every subtraction,
scalar input, fixed-point product and register remains in its guarded
domain. In particular, the scalar error cannot silently replace its
input-perturbation charge.

The represented upper endpoint is computed by adding the same represented
nonnegative $\log T$ register to the represented lower endpoint. Both
payoffs use the same normal draw, offset, constants and scalar table.
The exponential table is monotone: rounded node integers are nondecreasing,
and floor interpolation within a cell is nondecreasing and does not exceed
the next node. Positive subtraction, clipping and positive scaling preserve
this order. Hence the represented payoff difference is nonnegative before
its defensive positive subtraction, provided the independently checked
overflow guards hold. The final clamp yields a probability in $[0,1]$.

\section{Inner Gaussian discretisation}

\begin{lemma}[Normal-grid Kolmogorov bound]\label{lem:raw-cdf}
Let $L>0$, $a=\Phi(L)-\Phi(-L)$ and $M=2^{b_N}$. Define the exact
conditional midpoint quantiles
\begin{equation}
 z_j=\Phi^{-1}\!\left(\Phi(-L)+a\frac{2j+1}{2M}\right),
 \qquad j=0,\ldots,M-1.
\end{equation}
Let $\widetilde z_j$ satisfy $|\widetilde z_j-z_j|\leq e_N$,
and let $\widetilde N$ choose these values uniformly. If
$\nu_0\geq\sup_x\phi(x)$, then
\begin{equation}
 \sup_x|\Prb\{\widetilde N\leq x\}-\Phi(x)|
 \leq D_N:=\frac{1-a}{2}+\frac{a}{2M}+\nu_0e_N.
\end{equation}
An outward lower enclosure of $a$ may replace $a$ on the right.
\end{lemma}
\begin{proof}
First consider exact $z_j$. Between consecutive quantiles, its empirical
CDF is $j/M$ for some $j$, whereas the normal CDF ranges between the
neighbouring midpoint probabilities
$\Phi(-L)+a(j-1/2)/M$ and $\Phi(-L)+a(j+1/2)/M$.
Since $\Phi(-L)=(1-a)/2$, its deviation from $j/M$ is bounded by
$|(1-a)(1/2-j/M)|+a/(2M)$, which is at most
$(1-a)/2+a/(2M)$. Below the first quantile and above the last one,
the same bound follows directly from their midpoint probabilities and
the normal tail limits. At atoms the one-sided limits give the same bound.
If $F_M$ is this empirical CDF and $\widetilde F_M$ the rounded CDF,
then $F_M(x-e_N)\leq\widetilde F_M(x)\leq F_M(x+e_N)$,
regardless of duplicate rounded quantiles. Since
$|\Phi(x\pm e_N)-\Phi(x)|\leq\nu_0e_N$, the stated estimate follows.
Finally, $(1-a)/2+a/(2M)$ is decreasing in $a$ for $M\geq1$;
an outward lower endpoint therefore preserves the upper bound.
\end{proof}

\begin{proposition}[Nonsmooth payoff-difference transfer]\label{prop:raw-tv}
For every $m\in\mathbb R$, $\Delta\geq0$ and $\lambda>0$,
\begin{equation}
 |\E h_{m,\Delta}(\widetilde N)-\E h_{m,\Delta}(N)|\leq2D_N.
\end{equation}
\end{proposition}
\begin{proof}
Both maps $n\mapsto f(m+\Delta+\lambda n)$ and
$n\mapsto f(m+\lambda n)$ are continuous and nondecreasing,
with limits zero and one. Each consequently defines a positive
Lebesgue--Stieltjes derivative measure of mass one. Their difference
$h_{m,\Delta}$ has signed derivative $dh$ with total variation
$|dh|(\mathbb R)\leq2$. It is continuous and tends to zero at
both infinities. For the signed probability measure
$d(\widetilde F_M-\Phi)$, integration by parts, first on finite
intervals and then by bounded convergence against the finite measure
$|dh|$, gives
\begin{equation}
 \E h(\widetilde N)-\E h(N)
   =-\int_{\mathbb R}(\widetilde F_M(x)-\Phi(x))\,dh(x).
\end{equation}
Boundary terms vanish because both CDF differences vanish at infinities
and $h$ is bounded. The absolute integral is bounded by
$D_N|dh|(\mathbb R)\leq2D_N$. This proof permits atoms, coincident
rounded quantiles and nondifferentiable payoff kinks. It uses no empirical
price value and no local Gaussian smoothness of the payoff.
\end{proof}

% Supporting TeX fragment, intended for the existing raw-payoff proof.
% This contains a checked derivation, not a claim of theoretical priority.
\section{A smoothing-first arithmetic certificate}
\label{sec:raw-smoothing-first}

The complete pointwise graph certificate remains necessary for unsigned
domains, scalar input domains and the absence of modular overflow. Its
pointwise output error need not be the sharpest certificate of the prepared
mean. I now derive a second bound for the same emitted arithmetic and take
the smaller of the two bounds. This changes the financial certification,
prior enlargement and estimator allowance; it does not change the circuit.
The following argument separates rounding which depends on the inner
normal label from rounding which depends only on the outer coordinate.
Consequently, it does not differentiate a finite atomic law or a payoff
across its kinks.

Write $s=2^q$ and $u_q=s^{-1}$. The nearest dyadic rounding convention is
\begin{equation}
 \mathcal R_q(v)=s^{-1}\lfloor sv+1/2\rfloor.
\end{equation}
All quantities in this section are real values decoded from the guarded
integer graph, rather than integers interpreted modulo a register width.
Let $y$ be any realised outer-grid coordinate. The exact targets are
$m=m_G(y)$ and $\Delta=\log T(y)\geq0$. The shifted affine node has ideal
value $M=m+B_G$, where $B_G$ is the exact offset inherited from the retained
graph. Its implemented value is $\widehat M$. The implemented nonnegative
logarithm node is $\widehat\Delta$. Let the retained node certificate give
the uniform bounds
\begin{equation}
 |\widehat M-M|\leq e_m,
 \qquad |\widehat\Delta-\Delta|\leq e_\Delta,
 \qquad \widehat\Delta\geq0.
\end{equation}
The words ``realised outer-grid coordinate'' are essential: Gaussian
outer-node rounding is already included in the inherited financial
transfer. The arithmetic certificate treats those realised coordinates
as its exact inputs and does not charge their rounding a second time.

Put
\begin{equation}
 C=J-B_G,\quad O=J+\lambda L,\quad
 \widehat C=\mathcal R_q(C),\quad
 \widehat O=\mathcal R_q(O),
\end{equation}
and define $\epsilon_C=\widehat C-C$ and
$\epsilon_O=\widehat O-O$. Their certified absolute bounds are
$e_C$ and $e_O$. These constants are independent of the inner normal
label. In this graph the lower-endpoint additions and the upper-endpoint
addition are exact integer additions. Their real decoded values therefore
agree with ordinary real addition, once the retained whole-graph overflow
and domain guards pass. No additional addition-rounding term is present.

Let $(I_\lambda,d_\lambda)$ be exactly the coefficient pair returned by
the retained fixed-point scaling emitter and set
\begin{equation}
 \widehat\lambda=I_\lambda 2^{-d_\lambda},
 \qquad \delta_\lambda=\widehat\lambda-\lambda,
 \qquad
 e_{\rm floor}=\begin{cases}u_q,&d_\lambda>0,\\0,&d_\lambda=0.\end{cases}
\end{equation}
For the realised quantile word $v=\widetilde N+L\in[0,2L]$, the scale
node decodes to
\begin{equation}
 \mathsf S_\lambda(v)
 =s^{-1}\left\lfloor (sv)I_\lambda 2^{-d_\lambda}\right\rfloor
 =\widehat\lambda v+\rho(v),
 \qquad |\rho(v)|\leq e_{\rm floor}.
\end{equation}
When $d_\lambda=0$, the product is an integer and $\rho(v)=0$ exactly.
The inequality with $u_q$ is outward and remains valid at dyadic endpoints.
Thus, with
\begin{equation}
 E_s=2L|\delta_\lambda|+e_{\rm floor},\qquad
 \widehat m=\widehat M-B_G+\epsilon_C-\epsilon_O,
\end{equation}
the two conceptual signed log arguments before positive subtraction are
exactly
\begin{align}
 \widehat X_G
 &=\widehat M+\widehat C+
     \mathsf S_\lambda(\widetilde N+L)-\widehat O
   =\widehat m+\lambda\widetilde N+\zeta(\widetilde N),\\
 \widehat X_A
 &=\widehat X_G+\widehat\Delta,
 \qquad
 \zeta(n)=\delta_\lambda(n+L)+\rho(n+L),
 \qquad |\zeta(n)|\leq E_s.
\end{align}
The signed expressions are mathematical identities used to identify the
target. The circuit itself computes their positive parts by guarded
positive subtraction; it does not store an unguarded signed logarithm.
In particular,
\begin{equation}
 |\widehat m-m|\leq E_m:=e_m+e_C+e_O.
\end{equation}

\begin{lemma}[Scalar implementation error after exact payoff clamps]
\label{lem:raw-payoff-scalar-error}
Let $b=1/64$ and let the implemented exponential map $\widehat E$ satisfy
\begin{equation}
 \sup_{t\in[0,b]\cap s^{-1}\mathbb Z}
 |\widehat E(t)-e^t|\leq e_{\exp}.
\end{equation}
With $K=100$ and $H=1$, the graph payoff map obeys
\begin{equation}
 |\widehat f(x)-f(x)|\leq100e_{\exp}
\end{equation}
for every guarded decoded log argument $x$.
\end{lemma}
\begin{proof}
The positive-part map and the clamps at $b$ and one are exact and
$1$-Lipschitz. Both constants $b$ and one are exactly representable at
$q=28$. The scalar graph's second input clamp at $b$ is idempotent and
exact. The operation which subtracts one positively is exact because
the constant one is exact and the positive-subtraction guards pass.
Scaling by the integer $100$ has $d=0$, so it is an exact integer product
without a fixed-point floor. The difference between implemented and exact
exponentials is therefore multiplied by at most $100$ and is not enlarged
by either following clamp. The clamp decomposition established earlier
identifies the exact composed payoff with $f(x)$, completing the proof.
\end{proof}

The scalar certificate used in this lemma is its
\texttt{output\_error\_upper}, comprising nodal rounding, smooth
interpolation and the interpolation floor. It excludes the propagated
argument-error charge. Argument perturbations are treated separately
below. This distinction is required to avoid both omitting a perturbation
and charging it twice. Allow distinct bounds $e_{\exp,G}$ and
$e_{\exp,A}$, although the actual graph uses the same table and hence the
same bound on both sides.

\begin{theorem}[Smoothing-first conditional mean error]
\label{thm:raw-smoothing-first}
Assume the whole-graph domain and overflow certificate passes, the
quantile CDF bound is $D_N$, the shifted-node and coefficient bounds above
hold, and $\nu_0\geq\sup_x\phi(x)$. Let $p(y)$ be the exact probability
of the prepared good flag conditional on the realised outer coordinate.
Then
\begin{align}
 |p(y)-R_1(y)|
 &\leq\eta_{\rm smooth},\\
 \eta_{\rm smooth}
 &=\min\left\{1,
     100(e_{\exp,G}+e_{\exp,A})+202E_s+2D_N
     +\frac{\nu_0}{\lambda}(2E_m+e_\Delta)\right\}.
\end{align}
If $\eta_{\rm point}$ is the independently valid pointwise graph/CDF
bound, the smaller bound
$\eta=\min\{\eta_{\rm point},\eta_{\rm smooth}\}$ is valid uniformly
over the same realised outer grid.
\end{theorem}
\begin{proof}
The payoff is absolutely continuous, with its almost-everywhere derivative
\begin{equation}
 f'(x)=100e^x\mathbf1_{(0,\ell)}(x),
 \qquad \ell=\log(101/100).
\end{equation}
Thus $f$ is globally $101$-Lipschitz. Replacing
$\widehat X_G$ and $\widehat X_A$ by
$\widehat m+\lambda\widetilde N$ and
$\widehat m+\widehat\Delta+\lambda\widetilde N$ changes the difference
of exact payoffs by at most $2\cdot101E_s=202E_s$ for each normal label.
The scalar lemma changes the difference by at most
$100(e_{\exp,G}+e_{\exp,A})$. The final positive subtraction and clamp
are the metric projection onto $[0,1]$; that projection is $1$-Lipschitz
and fixes the exact payoff difference because
$\widehat\Delta\geq0$. They cannot enlarge this bound.

For completeness, the represented output is in fact the unprojected
difference of represented payoffs. The upper represented endpoint is
the lower endpoint plus the nonnegative word $\widehat\Delta$, without
overflow. Each preceding clamp and positive subtraction is monotone.
The rounded exponential node integers are nondecreasing. Within a cell,
the interpolator returns its base plus the floor of a nonnegative
increment. It is nondecreasing and cannot exceed the next node; at the
next cell it begins at that node. Hence the interpolator is monotone on
the complete dyadic input domain. Positive scaling and the final payoff
clamp preserve order, giving $0\leq\widehat f(\widehat X_G)
\leq\widehat f(\widehat X_A)\leq1$. Exact final subtraction therefore
produces their difference. The uniform comparator implements that
dyadic difference as a flag probability exactly.

After these pointwise replacements, the finite-inner-law target is
$\mathbb E h_{\widehat m,\widehat\Delta}(\widetilde N)$.
Both $\widehat m$ and $\widehat\Delta$ depend only on outer labels,
constants and outer arithmetic. They are independent of the inner normal
label, as is required here. The existing bounded-variation CDF lemma
applies at these represented endpoints, including atomic or duplicate
quantiles, and bounds its difference from
$\mathbb E h_{\widehat m,\widehat\Delta}(N)$ by $2D_N$.
Normal-node rounding is entirely charged by $D_N$ and does not enter
$E_s$ or $E_m$.

It remains to compare two continuous conditional means. Define
\begin{equation}
 P_\lambda(t)=\mathbb E f(t+\lambda N).
\end{equation}
The nonnegative derivative $f'$ has mass
$\int f'(x)\,dx=100(e^\ell-1)=1$. Convolution with the Gaussian density
therefore gives
\begin{equation}
 P_\lambda'(t)
 =\int_{\mathbb R} f'(x)\frac1\lambda
             \phi\left(\frac{x-t}{\lambda}\right)\,dx,
 \qquad 0\leq P_\lambda'(t)\leq\frac{\nu_0}{\lambda}.
\end{equation}
One justification of the derivative identity is to write the difference
quotient of $f(t+\lambda N)$ as an integral of its bounded almost-everywhere
derivative and use dominated convergence. Changing variables yields the
displayed convolution. Alternatively, the finite positive derivative
measure $df$ can be convolved with the smooth Gaussian. Either argument
integrates the continuous normal before applying the slope bound and
never differentiates the finite atomic normal approximation.

The Lipschitz bound for $P_\lambda$ now yields
\begin{align}
 &\left|P_\lambda(\widehat m+\widehat\Delta)
          -P_\lambda(\widehat m)
          -P_\lambda(m+\Delta)+P_\lambda(m)\right|\\
 &\qquad\leq\frac{\nu_0}{\lambda}
        \left(|\widehat m-m+\widehat\Delta-\Delta|
              +|\widehat m-m|\right)
 \leq\frac{\nu_0}{\lambda}(2E_m+e_\Delta).
\end{align}
Adding the three independently justified errors proves the untruncated
bound. Both the prepared probability and $R_1(y)$ lie in $[0,1]$, so
truncating the error at one is valid. Taking the smaller of two valid
uniform upper bounds is valid as well.
\end{proof}

The analytic financial coefficient, outer finite law, its external mass
and continuous financial transfer remain unchanged. The complete monetary
certificate is rebuilt with this new $\eta$, and the same ideal finite-mean
prior is enlarged by this new $\eta$. The phase-estimation allowance,
coverage plan and full logical resource charge must be rebuilt afterwards.
Neither a new estimator cost nor a strengthened comparative conclusion
follows merely from deriving a smaller arithmetic bound. This is a
project-specific derivation from established smoothing, monotonicity,
fixed-point rounding and CDF transfer principles; theoretical priority is
not asserted.

\section{Complete monetary certificate and quantum preparation}

Let $e_{\rm ar}$ be the complete graph's pointwise output-error upper
bound relative to the exact payoff difference at the realised outer
coordinates and the realised rounded normal draw. Normal-node rounding
is charged in Lemma~\ref{lem:raw-cdf}, while the retained graph treats
its realised quantile integers as exact inputs. Therefore
\begin{equation}
 \eta_{\rm point}=\min\{1,e_{\rm ar}+2D_N\}
\end{equation}
is the original valid pointwise bound. I use
$\eta=\min\{\eta_{\rm point},\eta_{\rm smooth}\}$
for the amended probability and monetary certificate. Both bounds apply
to the same ideal finite-outer-grid conditional mean $R_1$. The factor two in the
normal discrepancy is a sufficient bounded-variation estimate, not a
claim that the discrepancy is attained.

For each graph the preparation first places its outer Gaussian labels,
inner Gaussian label and $q$ comparator bits in a uniform superposition.
It computes the guarded numerical probability, XORs the comparison
$U_q<p_q$ into a separate good flag and reverses the complete arithmetic
body. Here $U_q$ is uniform on $\{0,\ldots,2^q-1\}$ and the
probability integer $p_q$ lies in $\{0,\ldots,2^q\}$, so the
conditional good probability is exactly $p_q/2^q$. The comparator's
endpoint-high-zero wire is allocated but excluded from the reflected
active core. The core consists of the actual uniform input wires,
comparator wires and actual good flag. All clean private arithmetic
wires may be reused for reflection scratch. Full controlled reflections,
the conditional global-minus gate, $A$ and $A^\dagger$, the synthesised
inverse Fourier transform, measurements and resets are included in
the complete canonical QPE charge. This retains the published amplitude
estimation and derivative-pricing framework \citep{brassard,rebentrost};
no new generic amplitude-estimation result is claimed.

Let the retained outer finite law have external factor in
$[a_-,a_+]$, midpoint $a_c$ and radius $a_r$. Let $\beta$ be its
price-independent complete continuous financial transfer and $e_0$
the analytic coefficient radius. The direct difference's normaliser
is $H=1$, and its whole deterministic undiscounted monetary error obeys
\begin{equation}
 D=e_0+\beta+a_+\eta+a_r.
\end{equation}
To derive this bound, let $r$ be the ideal finite-outer-grid mean correction,
$p$ the implemented flag probability and $a$ the exact external outer mass.
By construction, $|V-(c_0+ar)|\leq\beta$,
$|\widehat c_0-c_0|\leq e_0$, $|p-r|\leq\eta$ and
$|a-a_c|\leq a_r$, with $a\leq a_+$. Clip any reported probability
to $[0,1]$, which cannot increase its error from $p\in[0,1]$.
On the estimator's coverage event $|\widehat p-p|\leq\varepsilon_p$,
the decoder $\widehat V=\widehat c_0+a_c\widehat p$ satisfies
\begin{equation}
 |\widehat V-V|
 \leq e_0+\beta+a|p-r|+a|\widehat p-p|
             +|a-a_c|\,|\widehat p|
 \leq D+a_+\varepsilon_p.
\end{equation}
Thus the probabilistic estimator allowance is added exactly once; none
of its confidence budget is spent again on the deterministic bounds.
The inner finite Gaussian law is normalised and requires no further
external mass multiplier. The inherited ideal finite-mean prior $[l,u]$
is enlarged to $[\max(0,l-\eta),\min(1,u+\eta)]$. Neither a
basket-price answer nor an earlier sampled estimate enters this prior.
If $D<1/100$, the canonical QPE target is
$\varepsilon_p=(1/100-D)/a_+$ with per-price failure allowance
$1/3600$. Its retained explicit amplitude bound, outward Fourier error,
exact binomial amplification and postprocessing allowance certify the
full monetary radius. The controller records any exhausted budget,
domain failure or production exception as a rejection; none is repaired
after inspection within this protocol.

\section{Registered pointwise comparison and interpretation}

The outer financial parameters, selected exact rational factor bases,
$L=6$, outer quantile width ten, outer scalar width twelve and $q=28$
are fixed. Before fresh costs I registered inner normal widths
$\{8,10,12\}$, narrow-exponential address widths $\{4,6,8,10\}$
and both coherent lookup architectures. The first registered run used the
safe exponential output interval $[1,2]$; I retained its complete evidence
and separately froze the exact $[1,64/63]$ range amendment before new costs.
The amendment changes only semantic range bounds and resulting integer
headroom, while retaining the target, scalar ordinates, precision pool and
confidence policy. Its complete resource record is retained in deterministic
gzip. These are disclosed sequential design choices, not a blind initial
registration. This gives twelve configurations
per financial case and forty-eight in total, with ninety-six complete
architectural resource plans. All forty-eight configurations were admitted.
Every selected minimum uses inner width ten and exponential width ten.
Thus a wider precision search could strengthen this rival; the registered
finite minimum is not a global optimisation claim.

Under the original pointwise certificate, its minimum complete logical T caps are
$23\,346\,683\,625$, $88\,548\,390\,850$,
$94\,437\,597\,250$ and $417\,229\,413\,325$ in the retained
case order. The fixed shared-threshold representation is respectively
$10.67$, $5.63$, $5.50$ and $51.27$ per cent cheaper. This new rival
is $1.86$ and $0.80$ per cent cheaper than the strengthened numerical
conditional-price rival in the first and third cases. It is $0.39$ and
$99.16$ per cent more expensive in the second and fourth cases. A strongest
registered-rival comparison therefore selects the less costly numerical or
raw-payoff rival independently in each case.
These are matched sufficient logical resource caps under synthetic
contracts and the registered finite policies. They are not quantum
price outcomes, physical runtimes, a general theorem of algorithmic
superiority or evidence of practical quantum advantage.
\section{Rebuilt smoothing-first comparison}

I separately register the smoothing certificate after deriving its full
graph-to-target proof and disclosing the known preceding outcomes. All
48 literal graphs, scalar tables, financial models and 96 coherent
profiles are unchanged. Every deterministic budget, enlarged ideal-mean
prior and complete canonical QPE plan is rebuilt. Independent mathematical
recomputation checks every coefficient, exponential error decomposition,
normal discrepancy and selected allowance. A separate own-formula audit
checks all 96 confidence and full-charge plans. Whole-preparation checks
cover the original winners and the newly selected eight-bit normal circuit.

The amended raw minima are $23\,346\,683\,625$,
$88\,548\,390\,850$, $94\,122\,973\,950$ and
$208\,512\,831\,750$ logical T gates. Their selected inner normal widths
are ten, ten, eight and ten; all exponential widths are ten. The third
case changes its winning circuit and the eight-asset case changes QPE
resolution from 1,024 to 512, without changing that preparation. All
selected plans use 25 repetitions. The strongest admitted rival is raw
payoff in the first, third and fourth cases and numerical conditional
pricing in the second. The fixed shared-threshold encoding then saves
$10.67$, $5.26$, $5.18$ and $2.49$ per cent, respectively. Exact integers
and rational ratios remain retained.

These narrower positive findings supersede the earlier pointwise-only
comparison for headline claims. Their magnitude depends on the chosen
sufficient certificates and finite precision policies; every exponential
winner remains at the pool's upper boundary. Further certified error
refinements, shared efficient arithmetic and wider precision optimisation
may strengthen either opponent. The result establishes neither generic
query superiority, physical quantum advantage nor theoretical priority.

\appendix
% Optional sharpened transfer fragment. It is distinct from frozen costs.
\section{Exact variation of the common-normal payoff difference}
\label{sec:raw-exact-variation}

\begin{proposition}[Exact variation]
\label{prop:raw-exact-variation}
For $\Delta\geq0$, $m\in\mathbb R$ and $\lambda>0$, let
$\ell=\log(101/100)$ and
$h(n)=f(m+\Delta+\lambda n)-f(m+\lambda n)$. Then
\begin{equation}
 \operatorname{TV}(h)
 =2\min\{1,101(1-e^{-\Delta})\}.
\end{equation}
Consequently, the normal CDF transfer obeys the sufficient conditional
bound
\begin{equation}
 |\mathbb E h(\widetilde N)-\mathbb E h(N)|
 \leq2D_N\min\{1,101(1-e^{-\Delta})\}.
\end{equation}
\end{proposition}
\begin{proof}
An increasing affine change of the independent variable preserves total
variation, so it suffices to study
$g(x)=f(x+\Delta)-f(x)$. If $\Delta=0$, the function is identically
zero. Suppose first that $0<\Delta<\ell$. Its explicit pieces are
\begin{equation}
 g(x)=\begin{cases}
 0,&x\leq-\Delta,\\
 100(e^{x+\Delta}-1),&-\Delta<x\leq0,\\
 100e^x(e^\Delta-1),&0<x<\ell-\Delta,\\
 101-100e^x,&\ell-\Delta\leq x<\ell,\\
 0,&x\geq\ell.
 \end{cases}
\end{equation}
The pieces agree at every boundary. The first nonzero piece is increasing;
the middle piece is increasing; the last nonzero piece is decreasing.
The maximum at $x=\ell-\Delta$ is
$101-100e^{\ell-\Delta}=101(1-e^{-\Delta})<1$.
Because the continuous function increases from zero to this maximum and
then decreases to zero, its variation is twice its maximum.

If $\Delta\geq\ell$, the nonzero pieces instead are
\begin{equation}
 g(x)=\begin{cases}
 100(e^{x+\Delta}-1),&-\Delta<x<\ell-\Delta,\\
 1,&\ell-\Delta\leq x\leq0,\\
 101-100e^x,&0<x<\ell.
 \end{cases}
\end{equation}
Outside their union the function is zero. The plateau has height one;
for $\Delta=\ell$ it reduces to a single point. The variation is two.
At $\Delta=\ell$, $101(1-e^{-\ell})=1$, so the two formulae agree.
The CDF estimate follows from the same signed Stieltjes integration by
parts used for the original variation upper bound, replacing its bound
two by the exact variation just derived.
\end{proof}

This sharpening is conditional on the endpoint gap. A uniform certificate
may replace $\Delta$ by a valid upper bound because
$\min\{1,101(1-e^{-\Delta})\}$ is nondecreasing. If that upper bound
exceeds $\ell$, the uniform estimate is still $2D_N$. For the
smoothing-first certificate, the gap in the CDF step is the represented
nonnegative $\widehat\Delta$, so a bound for the true gap alone does not
justify a tighter represented-gap transfer without its arithmetic
enlargement. This optional sharpening has not been used in the frozen
pointwise resource results. Any resource claim which uses it requires a
separately recorded financial certificate, prior and estimator rebuild.

% Project-specific derivation for the retained literal compiler.
% This fragment makes no claim of generic quantum lower-bound novelty.
\section{Excluding still wider addresses within the fixed compiler}
\label{sec:raw-wider-address-exclusion}

I distinguish an exclusion within a fixed emitted implementation from an
algorithmic complexity lower bound. In this section the classical financial
models, $q=28$, narrow exponential span $b_0=1/64$, both payoff graphs,
normal widths $\{8,10,12\}$, outer precisions, two coherent unary-lookup
architectures and canonical-QPE controller are fixed. The scalar address
width is denoted by $b_e$, rather than by $b_0$.

The dyadic input interval contains $2^{q-6}+1$ representable arguments.
An address width $b_e$ gives $N=2^{b_e}$ cells, with remainder width
$q-6-b_e=22-b_e$. Hence this interpolator requires $1\leq b_e\leq22$.
The widths $11,12,13,14$ leave respectively $11,10,9,8$ remainder bits;
no new financial or smoothing assumption follows from this extension.
The fresh scalar tables still require checked rounding, monotonicity,
input domains and whole-graph absence of modular overflow.

\begin{lemma}[Literal selector count]
\label{lem:raw-selector-count}
Suppose a unary lookup has $b_e+1$ address bits and exactly the first
$N+1$ words are nonzero, where $N=2^{b_e}$ and $b_e\geq1$.
The retained recursively pruned selector tree visits exactly $N+b_e$
internal nodes. Its logical T count is
\begin{equation}
 \tau_{\rm look}(b_e)=
 \begin{cases}
 14(N+b_e),&\text{exact seven-T Toffoli architecture},\\
 8(N+b_e),&\text{unitary eight-T selector-pair architecture}.
 \end{cases}
\end{equation}
\end{lemma}
\begin{proof}
The root partitions the $2N$ address values into two halves. Its left
subtree contains all $N$ nonzero words and is a full binary tree with
$N-1$ internal nodes. The right subtree contains the single nonzero word
at address $N$ and consequently visits a path with $b_e$ internal nodes.
The root itself contributes one, giving $1+(N-1)+b_e=N+b_e$.
Each visited node has exactly one selector compute and one selector
erase. In the first architecture these are two seven-T Toffolis. In the
second they are a four-T clean-target compute and its four-T unitary
inverse. All changes of branch, word writes and root operations are
Clifford gates. The stated counts therefore follow from the literal
retained sequence, with no measurement or free lookup assumption.
\end{proof}

All packed narrow-exponential words are nonzero because their base
ordinates are at least the exact integer representing one. The final
endpoint word is also present; it has zero interpolation delta and its
nonzero base. Thus Lemma~\ref{lem:raw-selector-count} applies.
Each exponential node contains two lookup invocations, as its private
arithmetic is cleaned. The two payoff nodes and the subsequent cleanup
of the complete preparation give eight such invocations in $A$.
Let $a^{\rm fixed}_{j,k}$ be the retained logical T count of all other
blocks and the comparator in case $j$, normal width $k$ and the stated
architecture, including their cleanup. Their arguments, word widths,
tables and emitted sequences are independent of $b_e$. Discarding every
other T gate inside the exponential nodes gives
\begin{equation}
 T_A(b_e)\geq a^{\rm fixed}_{j,k}+c(2^{b_e}+b_e),
 \qquad c=112\text{ or }64.
\end{equation}
This is valid even when a wider address shortens its interpolation
multiplier, because that multiplier has been discarded completely from
the lower bound. It relies on this literal graph, without a hypothetical
optimisation which shares the two scalar evaluations.

Let $A_j=a_{+,j}$ be the unchanged outer mass upper bound, $e_{0,j}$
the coefficient radius, $\beta_j$ the financial-transfer upper bound,
$a_{r,j}$ the outer mass radius and $[l_j,u_j]$ the retained ideal
finite-mean interval. Give a competing wider implementation the
optimistic error $\eta=0$, and define
\begin{equation}
 d_{0,j}=e_{0,j}+\beta_j+a_{r,j},\qquad
 \varepsilon_{0,j}=\frac{1/100-d_{0,j}}{A_j}.
\end{equation}
Every admitted implementation has deterministic error at least $d_{0,j}$
and hence probability tolerance at most $\varepsilon_{0,j}$. Its enlarged
probability upper prior is at least $u_j$. The controller's variance
upper rule
\begin{equation}
 v(u)=\begin{cases}u(1-u),&0\leq u\leq1/2,\\1/4,&1/2<u\leq1\end{cases}
\end{equation}
is nondecreasing, as is its outward dyadic square-root rule. Therefore
let $M_{0,j}$ be the first supported power of two satisfying the retained
canonical-QPE sufficient amplitude bound at $(\varepsilon_{0,j},u_j)$.
Every active admitted plan has $M\geq M_{0,j}$. The retained inequalities
$A_j(u_j-l_j)/2>1/100-d_{0,j}$ and $u_j>2\varepsilon_{0,j}$ exclude
both controller forms of a zero-shot prior decision. These are checked
exactly for each of the four cases.

For confidence amplification, set $\pi_+=355/113$ and give the inverse
Fourier transform zero operator error. The controller's maximal certified
single-shot success probability is then $s_0=8/\pi_+^2$. Exact binomial
calculation shows that the first odd repeat count with majority failure
at most $1/3600$ is $r_0=25$. For a smaller success probability the
majority-failure probability is nonincreasing in the success parameter
and is therefore no better. This monotonicity follows, for example, by
coupling each success indicator as $\mathbf1\{U_i\leq s\}$ with common
independent uniform $U_i$. Every active plan thus has at least $r_0$
repetitions. This is a constraint of this certification controller;
it is not a lower bound on all conceivable amplitude estimators.

One canonical-QPE repetition calls $A$ once initially and calls both
$A$ and $A^\dagger$ in each of its $M-1$ controlled Grover steps.
Ignoring all reflection and Fourier T gates gives the complete bound
\begin{equation}
 T_{\rm complete}(b_e)
 \geq r_0(2M_{0,j}-1)
       \left(\min_{k\in\{8,10,12\}}a^{\rm fixed}_{j,k}
             +c(2^{b_e}+b_e)\right).
\end{equation}
This lower bound concerns the fully specified logical resource policy;
it does not convert sufficient monetary-error bounds into an algorithmic
quantum query lower bound.

The exact retained fields give $M_{0,j}=256,512,512,512$ in case order.
For the less expensive unitary architecture the fixed-block minima are
$1\,494\,706$, $3\,129\,342$, $3\,348\,190$ and $7\,808\,330$ T gates,
all attained at normal width eight. At address width fifteen the lower
complete counts are respectively
\begin{equation}
45\,898\,249\,950,\quad133\,692\,136\,050,\quad
139\,289\,173\,650,\quad253\,357\,254\,150.
\end{equation}
Each already exceeds an admitted earlier raw-payoff cap in its own case.
The exact seven-T architecture has still larger corresponding bounds.
Since $2^{b_e}+b_e$ is strictly increasing, all feasible widths fifteen
through twenty-two are excluded from improving those incumbents under
the stated fixed choices. Widths beyond twenty-two are outside this
fixed-point interpolator's input-resolution domain.

This exclusion does not cover omitted lower address widths
$\{1,2,3,5,7,9\}$, unequal widths for the two payoff nodes, other normal
widths, shared arithmetic, different word precisions, other lookup
architectures or other mean estimators. It can resolve an upper-boundary
concern for the registered extension without supporting a claim of a
global precision optimum, generic quantum superiority, theoretical
priority or practical quantum advantage.

% Checked supporting fragment. This replaces any identical-map ordering
% assertion for independently addressed PG and PA exponential tables.
% It is a project-specific derivation, not a claim of theoretical priority.
\section{Independent payoff-table precision and projected encoding}
\label{sec:raw-asymmetric-projection}

I allow the two exponential look-up tables in the common-normal payoff
difference to use different address widths. The two represented maps can
be individually monotone without preserving their mutual order. The
financial target remains the original no-arbitrage conditional correction;
the following projection makes the implemented asymmetric difference a
valid probability without assuming an ordering that has not been proved.

Throughout this section, $K=100$, $H=1$, $\lambda>0$, $q=28$ and
$s=2^q$. Define
\begin{equation}
 f(x)=\min\{1,100(e^x-1)_+\},\qquad
 \Pi(t)=\min\{1,\max\{0,t\}\},\qquad
 \ell=\log(101/100).
\end{equation}
Since $\ell<1/64$, the exact bounded-argument decomposition is
\begin{equation}
 f(x)=\min\{1,100[\exp(\min\{1/64,x_+\})-1]_+\}.
\end{equation}
Indeed, the two expressions agree for $x\leq0$, for
$0<x<\ell$, and for $x\geq\ell$. The last case follows because
the clamped argument is at least $\ell$ whenever the unclamped
argument is at least $\ell$. Consequently, the exponential only needs
its certified domain $[0,1/64]$, with
$1\leq e^t\leq64/63$ there, the last inequality following from
$e^t\leq(1-t)^{-1}$ for $0\leq t<1$.

\begin{lemma}[Exact projected integer difference]
\label{lem:raw-asymmetric-projection}
Suppose every unsigned operand and intermediate product satisfies the
retained domain and overflow certificate. Let $a,b$ be the represented
payoff words, decoded in $s^{-1}\mathbb Z\cap[0,1]$. Then the emitted
operation \texttt{clamp(positive\_sub($a$,$b$),1)} decodes exactly to
$\Pi(a-b)$. In particular, its range is $[0,1]$ even if $a<b$.
The uniform $q$-bit comparator encodes this value as a flag probability
exactly.
\end{lemma}
\begin{proof}
Let the word width be $w$ and write $A=sa$, $B=sb$. The subtraction
primitive produces $A-B$ modulo $2^w$ in temporary workspace. The
comparison flag is one exactly when $B<A$. Controlled copies of the
temporary difference to a zero output word therefore return $A-B$ when
$A>B$, and zero when $A\leq B$. When $A>B$, the retained unsigned
bounds ensure $0<A-B<2^w$, so the copied word is the ordinary integer
difference. When $A<B$, the modular subtraction is never copied and is
uncomputed with the comparison workspace. The temporary modular word
does not become a financial value. The subsequent clamp at the exactly
representable integer $s$ gives $\min\{s,\max\{0,A-B\}\}$.
The guarded operands are nonnegative and fit their registers; the
retained pre-mask magnitude guard also excludes an ambiguous overflow.
Thus decoding gives exactly $\Pi(a-b)$. For an output integer
$C\in\{0,\ldots,s\}$, exactly $C$ labels in the uniform set
$\{0,\ldots,s-1\}$ satisfy $U<C$. The good-flag probability is
$C/s$, including both endpoints. The clean inverse restores the
workspace after the flag is copied.
\end{proof}

The projection is nondecreasing and $1$-Lipschitz. For example, if
$u\leq v$, then its piecewise-linear definition gives
$0\leq\Pi(v)-\Pi(u)\leq v-u$. Reversing the arguments gives the
absolute-value inequality. For every $\Delta\geq0$,
$h_{m,\Delta}(n)=f(m+\Delta+\lambda n)-f(m+\lambda n)$ lies in
$[0,1]$, so $\Pi(h_{m,\Delta}(n))=h_{m,\Delta}(n)$.

\begin{lemma}[Separate scalar errors survive projection]
\label{lem:raw-asymmetric-scalar}
Let $\widehat E_G$ and $\widehat E_A$ be the two guarded dyadic
exponential implementations, with independently certified bounds
\begin{equation}
 \sup_{t\in[0,1/64]\cap s^{-1}\mathbb Z}
 |\widehat E_i(t)-e^t|\leq e_{\exp,i},\qquad i\in\{G,A\}.
\end{equation}
Let $\widehat f_i$ denote the exact emitted clamps and positive
subtractions around $\widehat E_i$. For any guarded decoded arguments
$x_G,x_A$,
\begin{equation}
 \left|\Pi(\widehat f_A(x_A)-\widehat f_G(x_G))
          -\Pi(f(x_A)-f(x_G))\right|
 \leq100(e_{\exp,G}+e_{\exp,A}).
\end{equation}
\end{lemma}
\begin{proof}
The constants zero, one and $1/64$ are exactly representable at
$q=28$. All payoff clamps and positive subtractions implement their
decoded real maps exactly, conditional on the whole-graph guards.
Integer scaling by $100$ has no fixed-point floor. Since these clamps
are $1$-Lipschitz, $|\widehat f_i(x)-f(x)|\leq100e_{\exp,i}$.
The triangle inequality for the difference and the $1$-Lipschitz
property of $\Pi$ give the claim. No comparison between
$\widehat f_G$ and $\widehat f_A$ is required. Each $e_{\exp,i}$
is the scalar's \texttt{output\_error\_upper}, comprising nodal
rounding, interpolation error and interpolation floor. Its propagated
argument-error charge is excluded here because that charge is treated
separately below.
\end{proof}

Fix a realised outer-grid coordinate $y$. Let $m=m_G(y)$,
$\Delta=\log T(y)\geq0$, and let the inherited shifted node have
ideal value $M=m+B_G$ and represented value $\widehat M$.
Write $\widehat\Delta$ for the represented logarithm. The guarded
graph gives
\begin{equation}
 |\widehat M-M|\leq e_m,\qquad
 |\widehat\Delta-\Delta|\leq e_\Delta,\qquad
 \widehat\Delta\geq0.
\end{equation}
These errors treat the realised outer coordinates as exact inputs.
Outer Gaussian-node rounding is already charged by the inherited
continuous-to-finite financial transfer and is not charged again here.
Let $C=J-B_G$, $O=J+\lambda L$ and let their represented constants
be $\widehat C$, $\widehat O$. If
$|\widehat C-C|\leq e_C$ and $|\widehat O-O|\leq e_O$, set
\begin{equation}
 \widehat m=\widehat M-B_G+(\widehat C-C)-(\widehat O-O),
 \qquad E_m=e_m+e_C+e_O.
\end{equation}
Thus $|\widehat m-m|\leq E_m$.

Let $(I_\lambda,d_\lambda)$ be the actual integer coefficient and
drop count of the scaling emitter, and put
\begin{equation}
 \widehat\lambda=I_\lambda2^{-d_\lambda},\qquad
 E_s=2L|\widehat\lambda-\lambda|
       +s^{-1}\mathbf1_{\{d_\lambda>0\}}.
\end{equation}
For a realised quantile $\widetilde N\in[-L,L]$, the exact decoded
scale node is
$\widehat\lambda(\widetilde N+L)+\rho(\widetilde N+L)$,
where $|\rho|\leq s^{-1}$ when $d_\lambda>0$, and $\rho=0$
otherwise. The guarded endpoint additions are exact integer additions.
Consequently their conceptual signed arguments are
\begin{align}
 \widehat X_G&=\widehat m+\lambda\widetilde N+
                    \zeta(\widetilde N),\qquad |\zeta|\leq E_s,\label{eq:raw-asymmetric-arguments}\\
 \widehat X_A&=\widehat X_G+\widehat\Delta.\nonumber
\end{align}
These are mathematical signed differences used to identify the financial
target. The circuit obtains their positive parts by the guarded unsigned
construction of Lemma~\ref{lem:raw-asymmetric-projection}; it does not
store an unguarded signed log argument.

\begin{theorem}[Conditional error for independently addressed tables]
\label{thm:raw-asymmetric-conditional}
Assume the whole-graph certificate passes and the realised finite normal
law has Kolmogorov discrepancy at most $D_N$ from the standard normal
law. Let $\nu_0\geq\sup_x\phi(x)$ and let $p(y)$ be the implemented
good-flag probability conditional on the realised outer coordinate. Then
\begin{equation}
 |p(y)-R_1(y)|\leq\eta_{\rm smooth}:=
 \min\left\{1,100(e_{\exp,G}+e_{\exp,A})+202E_s+2D_N+
       \frac{\nu_0}{\lambda}(2E_m+e_\Delta)\right\}.
\end{equation}
The minimum of this bound and an independently valid pointwise
graph/CDF bound is a valid uniform conditional certificate.
\end{theorem}
\begin{proof}
The payoff is absolutely continuous, with almost-everywhere derivative
$f'(x)=100e^x\mathbf1_{(0,\ell)}(x)$, and is globally
$101$-Lipschitz. Replacing the two arguments in
\eqref{eq:raw-asymmetric-arguments} by
$\widehat m+\lambda\widetilde N$ and
$\widehat m+\widehat\Delta+\lambda\widetilde N$ changes the exact
payoff difference by at most $202E_s$ per normal label. Projection
cannot enlarge this error and fixes the latter difference because
$\widehat\Delta\geq0$. Lemma~\ref{lem:raw-asymmetric-scalar}
adds at most $100(e_{\exp,G}+e_{\exp,A})$.

The remaining finite-inner-law target is
$\mathbb E h_{\widehat m,\widehat\Delta}(\widetilde N)$. The
parameters $\widehat m$ and $\widehat\Delta$ depend only on outer
labels, not on the inner normal label. Each summand of $h$ is
continuous and nondecreasing with total variation one, so
$\operatorname{TV}(h)\leq2$. Stieltjes integration by parts gives
\begin{equation}
 \mathbb E h(\widetilde N)-\mathbb E h(N)
   =-\int_{\mathbb R}(F_{\widetilde N}(x)-\Phi(x))\,dh(x).
\end{equation}
The boundary term vanishes because $h$ is bounded and both CDFs tend
to the same limits at infinity. The continuous function $h$ assigns
no atom to its derivative measure, so the identity accommodates
atoms of the finite law and duplicate quantiles. Its absolute value
is at most $2D_N$. Normal rounding is charged solely within $D_N$.

Define $P_\lambda(t)=\mathbb E f(t+\lambda N)$. The nonnegative
measure $df$ has mass
$100(e^\ell-1)=1$. Gaussian convolution gives
\begin{equation}
 P_\lambda'(t)=\int_{\mathbb R}
       \frac1\lambda\phi\left(\frac{x-t}{\lambda}\right)df(x),
 \qquad 0\leq P_\lambda'(t)\leq\frac{\nu_0}{\lambda}.
\end{equation}
Differentiation is justified by the bounded smooth Gaussian kernel
against the finite measure $df$; equivalently, one may use the
bounded derivative of $f$ and dominated convergence. Thus the two
continuous endpoint perturbations contribute at most
$\nu_0(2E_m+e_\Delta)/\lambda$. Adding the independently justified
terms proves the uncapped bound. Both $p(y)$ and $R_1(y)$ lie in
$[0,1]$, so capping the error at one is valid. The minimum of two
valid uniform upper bounds remains valid.
\end{proof}

\begin{proposition}[Prospective common-shift refinement]
\label{prop:raw-common-shift-refinement}
For $R(m,\Delta)=P_\lambda(m+\Delta)-P_\lambda(m)$ with
$\Delta\geq0$, the continuous outer perturbation above also satisfies
\begin{equation}
 |R(\widehat m,\widehat\Delta)-R(m,\Delta)|
 \leq\frac{\nu_0}{\lambda}(E_m+e_\Delta).
\end{equation}
If $\nu_1\geq\sup_x|\phi'(x)|$ and a certified uniform bound
$\Delta\leq\Delta_*$ is available, the common-shift coefficient
can be replaced by
\begin{equation}
 \min\left\{\frac{\nu_0}{\lambda},
              \frac{\nu_1\Delta_*}{\lambda^2}\right\}E_m
       +\frac{\nu_0}{\lambda}e_\Delta.
\end{equation}
\end{proposition}
\begin{proof}
The previous convolution identity bounds $P_\lambda'$ in
$[0,\nu_0/\lambda]$. Hence
$|\partial_mR|=|P_\lambda'(m+\Delta)-P_\lambda'(m)|
\leq\nu_0/\lambda$, whereas
$|\partial_\Delta R|=P_\lambda'(m+\Delta)\leq\nu_0/\lambda$.
Split the perturbation into a gap change at $\widehat m$ followed by
a common-shift change at the true nonnegative gap $\Delta$, and
integrate these derivative bounds along the two intervals. Furthermore,
differentiating the Gaussian convolution against the probability
measure $df$ gives $|P_\lambda''|\leq\nu_1/\lambda^2$.
Integrating $P_\lambda''$ between $m$ and $m+\Delta$ bounds
$|\partial_mR|$ also by $\nu_1\Delta/\lambda^2$.
The uniform upper bound $\Delta_*$ gives the refined coefficient.
\end{proof}

This refinement is an analytic error bound, not an automatic improvement
in the full estimation plan. It must be selected and recorded in a
protocol before any resource outcomes using it are generated. The
financial coefficient, external outer-law mass, continuous transfer,
mean prior, estimator accuracy, confidence and full clean circuit costs
must be retained or rebuilt consistently. The projection argument
justifies asymmetric scalar precision without claiming cancellation of
the two distinct table errors, a new amplitude-estimation query law, or
a physical quantum advantage.

\section{Registered independent-width comparison}

The prospective common-shift refinement is not selected in this phase.
I retain the conservative $2E_m$ policy across the previously matched
opponents and register independent $b_G,b_A\in\{8,9,10,11,12\}$ with
normal widths $\{8,10,12\}$. Forty-eight diagonal settings reference exact
earlier evidence; the other 252 are emitted freshly. All 300 settings
are admitted, with both architectures and every pointwise/smoothing plan
retained. The union with the preceding 96 common-width settings has 348
distinct precision configurations.

Every new minimum uses $b_G=b_A=9$; the normal widths are ten, ten,
eight and ten. Thus the lower cap follows from the previously omitted
symmetric nine-bit setting, rather than a demonstrated advantage from
unequal allocation. Complete raw caps in case order are
$23\,240\,600\,025$, $88\,336\,016\,050$,
$93\,933\,514\,350$ and $208\,323\,372\,150$.
The numerical rival remains cheaper in the second case. Under the historical additive clean templates, against the
cheaper admitted rival across these studies, shared-threshold savings
are $10.26$, $5.26$, $4.99$ and $2.40$ per cent. All selected resolutions
and 25 confidence repetitions are unchanged.

The exact scalar witness at $t=43691/268435456$ gives exponential
integers 268479151 at width eight and 268479150 at width twelve.
After subtracting one and scaling by 100, an upper-side width-twelve
payoff and lower-side width-eight payoff at equal arguments therefore
give an unprojected difference of $-100$ integer units. Final projection
returns zero exactly. This proves a scalar-domain ordering failure,
not its occurrence in the retained financial grids. Independent literal
projection/comparator/inverse replay covers 192 basis assignments,
including dirty good flags and both architectures, with clean workspace.

Independent mathematics recomputes all 252 new conditional bounds,
37,899 scalar/quantile nodes and 10,752 closed/left CDF endpoints.
A separate resource replay reconstructs 504 profiles, 1,008 complete
QPE charges and all eight winning clean preparations. These are sufficient
logical implementation comparisons with complete monetary guarantees,
not hardware observations, a global precision optimum or a priority claim.

% Prospective resource exclusion, separate from registered numerical outcomes.
% Applies only to the retained literal graph/compiler/certification policy.
\section{A one-sided exclusion for independently addressed payoff tables}
\label{sec:raw-asymmetric-high-width-exclusion}

I extend the previous common-width exclusion to independently addressed
payoff exponential nodes. This argument is a lower bound on the logical
T count charged by the specified compiler and certification controller.
It is not a lower bound on derivative-pricing algorithms or amplitude
estimation in general. The four financial models, chosen factor bases,
$q=28$, outer precisions, inner normal widths $\{8,10,12\}$, complete
raw-payoff graph, clean unary-look-up architectures and canonical-QPE
controller are fixed. Only the two narrow exponential address widths,
denoted by $b_G$ and $b_A$, are allowed to vary independently.

\begin{lemma}[Width and non-exponential invariance]
\label{lem:raw-asymmetric-fixed-width}
Within this graph, changing either narrow exponential address width in
$\{1,\ldots,22\}$ leaves all ideal node ranges, all global arithmetic
word widths, all other scalar tables, every non-exponential emitted block
and the final $q$-bit flag comparator unchanged at fixed financial case
and inner normal width.
\end{lemma}
\begin{proof}
The narrow exponential input is exactly clamped to $[0,1/64]$. Its
dyadic arguments have integer words $t\in\{0,\ldots,2^{22}\}$.
An address width $b$ has $2^b$ interpolation cells and remainder width
$28-6-b=22-b$. The emitted scalar takes the address slice beginning
at bit $22-b$ and ending at bit $22$, inclusive. This is a feasible
slice of every retained word precisely for $1\leq b\leq22$.
The ideal exponential output range remains $[1,64/63]$, irrespective
of $b$. All subsequent ideal payoff ranges and their arithmetic
operations are therefore unchanged.

The graph finaliser determines its global headroom from the largest
absolute ideal range endpoint, with a fixed safety factor, and not from
scalar address widths. Every retained node is assigned this same global
word width. Consequently, changing $b_G$ or $b_A$ cannot change that
width or any non-exponential node's operand widths, alias pattern,
operation, constants or table. The two variable scalar nodes still
return words of that fixed width. Their privately allocated temporary
workspace is cleaned before another node is evaluated; it cannot alter
the arithmetic operand widths of another node. The complete preparation
is reconstructed as the unchanged forward non-exponential blocks and
their inverses, together with the variable scalar blocks and the same
final comparator. Its comparator operates on $q+1$ wires per operand,
with $q$ fixed, so its T charge is also unchanged.

The independent evidence check additionally compares the full retained
non-exponential graph and block records across all 504 fresh architecture
profiles and 96 referenced diagonal profiles. At each case/normal-width/
architecture combination their hashes and counts coincide exactly.
These comparisons corroborate the source-level invariance; they do not
replace its argument for omitted address widths.
\end{proof}

\begin{lemma}[Per-side selector charge]
\label{lem:raw-asymmetric-one-side-selector}
For a narrow exponential address width $b\in\{1,\ldots,22\}$,
the retained complete preparation $A$ contains four unary look-ups for
that side. Discarding all other operations inside that scalar node, its
selector T contribution is
\begin{equation}
 c_{\rm side}(2^b+b),\qquad
 c_{\rm side}=\begin{cases}
 56,&\text{exact seven-T Toffoli architecture},\\
 32,&\text{unitary eight-T selector-pair architecture}.
 \end{cases}
\end{equation}
\end{lemma}
\begin{proof}
Each packed table word contains a base ordinate equal to the nearest
nonnegative dyadic integer representing $e^x$, for $0\leq x\leq1/64$.
The base is at least $2^{28}$, so every one of the first $2^b+1$
packed words is nonzero. The last endpoint word has a zero interpolation
delta but its base remains nonzero. Thus the pruned selector tree sees
the complete left half with $2^b$ leaves and one nonzero leaf in the
right half. The left half contributes $2^b-1$ internal nodes; the
right path contributes $b$; the root contributes one. The total is
$2^b+b$. Each internal selector compute/erase pair costs fourteen T
gates in the exact-Toffoli architecture or eight T gates in the coherent
unitary architecture.

The scalar emitter performs a look-up for its packed interpolation data
and uncomputes that look-up after its output has been copied. There
are therefore two invocations per forward scalar node, also recorded
in every retained block profile. Cleanup of the complete preparation
includes the inverse scalar node and doubles this count to four.
Multiplying fourteen or eight by four gives the stated coefficients.
Word writes and branch changes are Clifford operations. Interpolation
and all other scalar arithmetic have been discarded, so reducing the
remainder width at larger $b$ cannot invalidate this lower bound.
\end{proof}

Let $a^{\rm fixed}_{j,k,\mathcal A}$ be the complete preparation's
T charge of every block other than the two narrow exponential nodes,
including its final comparator, for financial case $j$, normal width
$k$ and architecture $\mathcal A$. This includes forward execution
and cleanup. For either side with address width $b$, the preceding
lemmas imply
\begin{equation}
 T_A(b_G,b_A)\geq a^{\rm fixed}_{j,k,\mathcal A}
                         +c_{\rm side,\mathcal A}(2^b+b).
\label{eq:raw-asymmetric-one-side-A-floor}
\end{equation}
This discards the other side completely. Hence it is valid for every
feasible other address width from one to twenty-two, including equal
or unequal widths and without assuming represented payoff ordering.

I next give the prospective wider implementation the optimistic uniform
arithmetic and inner-law error $\eta=0$. Let $A_j$ be the unchanged
upper outer-law mass factor, $e_{0,j}$ the coefficient radius,
$\beta_j$ the financial-transfer allowance, $a_{r,j}$ the outer
mass radius and $[l_j,u_j]$ the unchanged ideal finite-mean prior.
Define
\begin{equation}
 d_{0,j}=e_{0,j}+\beta_j+a_{r,j},\qquad
 \varepsilon_{0,j}=\frac{1/100-d_{0,j}}{A_j}.
\end{equation}
Every admitted graph has deterministic monetary allowance at least
$d_{0,j}$ and probability tolerance at most $\varepsilon_{0,j}$.
Its implemented prior upper endpoint is at least $u_j$, while its
implemented prior interval is at least as wide as $[l_j,u_j]$.
The retained inequalities
\begin{equation}
 \frac{A_j(u_j-l_j)}2>\frac1{100}-d_{0,j},\qquad
 u_j>2\varepsilon_{0,j}
\label{eq:raw-asymmetric-zero-shot-exclusions}
\end{equation}
hold exactly for all four cases. The first excludes the two-sided
prior-midpoint controller decision, and the second excludes the
one-sided prior-midpoint decision. No zero-shot price decision within
this retained policy can evade the active-estimation charge below.

Put $\pi_+=355/113$ and
\begin{equation}
 v(u)=\begin{cases}u(1-u),&u\leq1/2,\\1/4,&u>1/2.\end{cases}
\end{equation}
The controller's outward dyadic square root of $v(u)$ is monotone in
$u$. Let $M_{0,j}$ be the first supported power of two satisfying
\begin{equation}
 \frac{2\pi_+\,\overline{\sqrt{v(u_j)}}}{M_{0,j}}
       +\frac{\pi_+^2}{M_{0,j}^2}+10^{-60}
       \leq\varepsilon_{0,j}.
\end{equation}
Because both the prior upper endpoint and the outward square root
cannot decrease, and the available tolerance cannot increase, every
admitted active plan has $M\geq M_{0,j}$. The independent exact
calculation gives $M_{0,j}=256,512,512,512$ in case order.

For repetitions, give the Fourier circuit the optimistic zero norm
error. The maximal certified single-shot success is
$s_0=8/\pi_+^2$. For every odd $r$, its majority-failure expression is
\begin{equation}
 B_r(s_0)=\sum_{k=(r+1)/2}^{r}\binom rk(1-s_0)^ks_0^{r-k}.
\end{equation}
Exact rational evaluation gives $B_r(s_0)>1/3600$ for all odd
$r\in\{1,3,\ldots,23\}$ and $B_{25}(s_0)\leq1/3600$.
Thus $r_0=25$ is the least admitted odd repeat count, even with perfect
Fourier synthesis. For any smaller success guarantee, coupling the
Bernoulli outcomes with common independent uniforms proves that
majority failure cannot improve. Every active plan therefore has
$r\geq25$. This is a floor imposed by the specified certification
controller, rather than by all quantum estimators.

\begin{theorem}[Either-side high-address exclusion]
\label{thm:raw-asymmetric-high-width-exclusion}
Under the fixed graph and controller above, if either $b_G$ or $b_A$
belongs to $\{15,\ldots,22\}$, no feasible choice of the other
address width in $\{1,\ldots,22\}$ improves the retained PG/PA
nine-bit incumbent in any of the four cases.
\end{theorem}
\begin{proof}
Each canonical-QPE repetition has one initial preparation call and
$M-1$ controlled Grover calls containing one $A$ and one
$A^\dagger$. Discarding all reflections, Fourier operations and
other T gates gives
\begin{equation}
 T_{\rm complete}(b_G,b_A)
 \geq25(2M_{0,j}-1)
    \left(\min_{k\in\{8,10,12\}}a^{\rm fixed}_{j,k,\mathcal A}
           +c_{\rm side,\mathcal A}(2^b+b)\right).
\label{eq:raw-asymmetric-complete-floor}
\end{equation}
The fixed counts for the exact-Toffoli architecture are no smaller
than those for the unitary architecture at each normal width; both
use identical arithmetic and the latter has the lower selector charge.
It therefore suffices to use the more optimistic unitary architecture
with coefficient thirty-two. Its fixed-block minima are
$1\,494\,706$, $3\,129\,342$, $3\,348\,190$ and
$7\,808\,330$, respectively, all attained at normal width eight.
At $b=15$, equation~\eqref{eq:raw-asymmetric-complete-floor} gives
the four respective T floors
\begin{equation}
 32\,496\,559\,550,\quad106\,862\,528\,850,\quad
 112\,459\,566\,450,\quad226\,527\,646\,950.
\end{equation}
The retained PG/PA nine-bit incumbents are respectively
\begin{equation}
 23\,240\,600\,025,\quad88\,336\,016\,050,\quad
 93\,933\,514\,350,\quad208\,323\,372\,150.
\end{equation}
Every floor strictly exceeds its incumbent. Since $2^b+b$ is strictly
increasing, the same holds for all feasible $b\geq15$. The other
side was discarded from the floor and remains unrestricted throughout
its feasible range. This proves the result for either side separately
and hence whenever at least one side has such an address width.
\end{proof}

This prospective exclusion does not evaluate new wider financial
certificates and does not modify any registered numerical outcome.
It leaves asymmetric widths thirteen and fourteen and omitted lower
widths unresolved. It also does not cover other normal widths,
word precisions, shared scalar arithmetic, non-unary look-ups, other
mean estimators, hardware implementations or optimisations that change
the retained graph. Consequently, it provides no global precision
optimum, generic quantum query lower bound, theoretical-priority claim
or practical quantum advantage.

\section{Common-information stability of a positive Gaussian correction}
\label{sec:common-information-stability}

I separate the mathematics of a common shift from the finite arithmetic
which implements it. A common shift is a perturbation of both endpoint
arguments by the same amount; it need not be probabilistically independent
of the outer financial coordinate or of the gap. The essential condition
for Gaussian smoothing is that, after conditioning on the outer labels
and any fixed strike label, the shift and gap are independent of the
inner normal label. A label-dependent scaling remainder does not satisfy
this condition and receives a separate allowance.

Let $\mu$ be a positive finite Borel measure of mass one, let $\lambda>0$,
and define
\begin{equation}
 P_{\lambda,\mu}(m)=\int_{\mathbb R}
       \Phi((m-a)/\lambda)\,\mu(da),\qquad
 R_{\lambda,\mu}(m,\Delta)=P_{\lambda,\mu}(m+\Delta)
                              -P_{\lambda,\mu}(m).
 \label{eq:common-information-definition}
\end{equation}
This includes a fixed strike, a finite positive strike mixture and a
bounded increasing payoff normalised to total variation one. In the last
case $\mu=df$, where $f(-\infty)=0$ and $f(+\infty)=1$; Tonelli's theorem
identifies $P_{\lambda,\mu}(m)$ with $\mathbb E f(m+\lambda N)$.
These identities and Gaussian convolution are established mechanisms.
The calculation below concerns their shared perturbations and their
specific guarded pricing implementations; it is not a claim of priority
for convolution or Lipschitz estimates.

Choose $\nu_0\geq\sup_x\phi(x)$ and
$\nu_1\geq\sup_x|\phi'(x)|$, and put
\begin{equation}
 \kappa_0=\frac{\nu_0}{\lambda},\qquad
 \kappa_1=\frac{\nu_1}{\lambda^2},\qquad
 C(d)=\min\{\kappa_0,\kappa_1d\}\quad(d\geq0).
 \label{eq:common-information-constants}
\end{equation}

\begin{theorem}[Common-shift and gap stability]
\label{thm:common-information-stability}
For finite $m,\widehat m,\Delta,\widehat\Delta$ with $\Delta\geq0$, let
$|\widehat m-m|\leq E_m$ and
$|\widehat\Delta-\Delta|\leq e_\Delta$. The represented gap may be
negative. Then
\begin{equation}
 |R_{\lambda,\mu}(\widehat m,\widehat\Delta)
                       -R_{\lambda,\mu}(m,\Delta)|
 \leq C(\Delta)E_m+\kappa_0e_\Delta.
 \label{eq:common-information-stability}
\end{equation}
Consequently a uniform true-gap bound $0\leq\Delta\leq\Delta_*$ permits
the coefficient $C(\Delta_*)$. If $\widehat\Delta\geq0$ as well, the
coefficient in \eqref{eq:common-information-stability} can be replaced
by $C(\min\{\Delta,\widehat\Delta\})$.
Let $\Pi_{[0,B]}$ be metric projection onto $[0,B]$, where $B>0$ and
$R_{\lambda,\mu}(m,\Delta)\in[0,B]$. The same upper bound applies after
projecting the represented difference onto $[0,B]$. Projection need not
be inserted into the implemented gap register.
\end{theorem}
\begin{proof}
The first and second derivatives of the Gaussian kernel are bounded
uniformly in $a$ by $\kappa_0$ and $\kappa_1$. Difference quotients of the kernel
are dominated by these constants by the mean value theorem. Since
$\mu$ is finite, dominated convergence justifies differentiation under
its integral twice, giving
\begin{align}
 P_{\lambda,\mu}'(m)&=\int\lambda^{-1}
                 \phi((m-a)/\lambda)\,\mu(da),\nonumber\\
 P_{\lambda,\mu}''(m)&=\int\lambda^{-2}
                 \phi'((m-a)/\lambda)\,\mu(da),\qquad
 0\leq P_{\lambda,\mu}'\leq \kappa_0,\quad
 |P_{\lambda,\mu}''|\leq \kappa_1.
 \label{eq:common-information-derivatives}
\end{align}
The difference of two numbers in $[0,\kappa_0]$ has absolute value at most
$\kappa_0$. The fundamental theorem of calculus also bounds the difference
of the derivatives at arguments separated by a nonnegative distance $d$
by $\kappa_1d$. Thus
\begin{equation}
 |\partial_m R_{\lambda,\mu}(m,d)|
   =|P_{\lambda,\mu}'(m+d)-P_{\lambda,\mu}'(m)|\leq C(d),\qquad
 |\partial_dR_{\lambda,\mu}(m,d)|\leq \kappa_0.
 \label{eq:common-information-partial-bounds}
\end{equation}
For the gap partial derivative the bound holds for all real $d$.
Split the difference exactly as
\begin{align}
 R(\widehat m,\widehat\Delta)-R(m,\Delta)
 &=R(\widehat m,\widehat\Delta)-R(\widehat m,\Delta)
       +R(\widehat m,\Delta)-R(m,\Delta).
 \label{eq:common-information-path}
\end{align}
Integrating the gap derivative along the first segment gives
$\kappa_0|\widehat\Delta-\Delta|$, including a segment crossing zero.
Integrating the common-shift derivative along the second segment, whose
gap is the true nonnegative $\Delta$, gives
$C(\Delta)|\widehat m-m|$. This proves the stated inequality. If the
represented gap is nonnegative, reversing the order of the two segments
gives $C(\widehat\Delta)|\widehat m-m|+\kappa_0|\widehat\Delta-\Delta|$;
the smaller valid bound gives the refinement. Since $C$ is nondecreasing,
a valid upper gap bound may replace the gap. Finally, metric projection
is $1$-Lipschitz and fixes the true correction. Applying it to the
represented difference cannot increase its distance from that correction.
No cancellation of two distinct scalar-table approximation errors is
asserted.
\end{proof}

\begin{lemma}[Common inner perturbations of capped payoffs]
\label{lem:common-information-inner-payoff}
For the normalised cap
$f(x)=H^{-1}\min\{K(e^x-1)_+,H\}$, put
$\ell=\log(1+H/K)$ and $L_f=(K+H)/H$. For every $d\geq0$, the map
$g_d(x)=f(x+d)-f(x)$ is globally $L_f$-Lipschitz. For any real
perturbation $\zeta$, which may depend on the inner label,
\begin{equation}
 |g_d(x+\zeta)-g_d(x)|\leq L_f|\zeta|.
 \label{eq:common-information-inner-payoff}
\end{equation}
\end{lemma}
\begin{proof}
The payoff is absolutely continuous and has almost-everywhere derivative
$f'(x)=(K/H)e^x\mathbf1_{(0,\ell)}(x)$, which lies in $[0,L_f]$.
The difference $g_d$ is absolutely continuous on every finite interval
and its almost-everywhere derivative is $f'(x+d)-f'(x)$.
Its absolute value is at most $L_f$, since both terms belong to the same
interval $[0,L_f]$. Integrating that derivative between $x$ and $x+\zeta$
proves the inequality. At $d=0$ the difference is identically zero.
This argument acts on the continuous payoff as a pointwise function and
does not differentiate a finite atomic probability law.
\end{proof}

\subsection{Matched finite raw-payoff implementation}

Retain $K=100$, $H=1$, $q=28$ and the guarded graph. Conditional on an
actual outer coordinate $y$, its decoded endpoints satisfy exactly
\begin{equation}
 \widehat X_G=\widehat m+\lambda\widetilde N+\zeta(\widetilde N),\qquad
 \widehat X_A=\widehat X_G+\widehat\Delta,\qquad
 |\widehat m-m|\leq E_m=e_m+e_C+e_O,\quad
 \widehat\Delta\geq0.
 \label{eq:common-information-raw-endpoints}
\end{equation}
Here $C=J-B_G$ and $O=J+\lambda L$ are constants, and their errors
$e_C,e_O$ are independent of the inner label. The exact emitted scaling
coefficient is $\widehat\lambda=I_\lambda2^{-d_\lambda}$ and
\begin{equation}
 |\zeta|\leq E_s=2L|\widehat\lambda-\lambda|
                    +2^{-q}\mathbf1_{\{d_\lambda>0\}}.
 \label{eq:common-information-scale-error}
\end{equation}
Let the two distinct scalar maps have intrinsic exponential errors
$e_{\exp,G},e_{\exp,A}$ on their guarded domains. Intrinsic errors exclude
propagated argument perturbations. Let the rounded inner normal law have
both right-CDF and left-CDF discrepancy at most $D_N$ from $\Phi$.
The prepared flag encodes the exact projected dyadic payoff difference.
Its conditional probability $p(y)$ obeys
\begin{equation}
 |p(y)-R_1(y)|\leq\eta_{\rm raw}:=
 \min\left\{1,100(e_{\exp,G}+e_{\exp,A})+101E_s+2D_N
                       +C(\Delta_*)E_m+\kappa_0e_\Delta\right\}.
 \label{eq:common-information-raw-bound}
\end{equation}
An independently valid pointwise or earlier smoothing bound may be
included in this minimum.

To prove \eqref{eq:common-information-raw-bound}, replace the two
represented scalar maps by exact payoffs at their represented arguments.
Their separate errors contribute $100(e_{\exp,G}+e_{\exp,A})$.
Lemma~\ref{lem:common-information-inner-payoff} removes the common
label-dependent $\zeta$ for a charge $101E_s$. Projection is
$1$-Lipschitz and fixes the resulting ideal nonnegative difference.
That difference as a function of the normal coordinate has bounded
variation at most two; integration by parts against the signed CDF
difference therefore charges at most $2D_N$. This statement retains
atomic label multiplicities and uses the left limit at strict endpoints.
Only after this finite-to-continuous transfer is the Gaussian
convolution differentiated. Theorem~\ref{thm:common-information-stability}
then charges $C(\Delta_*)E_m+\kappa_0e_\Delta$. The total is capped at one
because both target and prepared probability belong to $[0,1]$.
Normal quantile rounding is contained in $D_N$; it is not also included
in $E_s$ or $E_m$. Outer quantile rounding remains solely in the existing
financial outer-law transfer.

\subsection{Matched finite numerical conditional-price implementation}

For the retained numerical pool, the cap and the positive monetary
normaliser are exactly $H=B_1=1$. Consequently its two scaling nodes
multiply by the exact integer one and introduce neither coefficient
rounding nor a fixed-point floor. Write the scalar domains as
$[-B_G,B_G]$ and $[-B_A,B_A]$. The encoded nodes satisfy
\begin{equation}
 \widehat M_A=\widehat M_G+\widehat C_{AG}+\widehat\Delta,\qquad
 C_{AG}=B_A-B_G,\qquad
 \widehat m=\widehat M_G-B_G,\qquad
 \widehat d=\widehat\Delta+(\widehat C_{AG}-C_{AG}).
 \label{eq:common-information-numerical-endpoints}
\end{equation}
Suppose the retained actual-range certificate additionally shows that
the input clamps preceding both price tables are inactive on every
actual guarded input. A function-input domain check after a clamp alone
does not establish this premise. Let $e_{P,G},e_{P,A}$ be the intrinsic
conditional-price table errors, $e_{AG}$ the constant error and $e_m$
the lower encoded node error. The prepared projected difference then
satisfies
\begin{equation}
 |p_{\rm num}(y)-R_1(y)|\leq\eta_{\rm num}:=
 \min\{1,e_{P,G}+e_{P,A}+C(\Delta_*)e_m
                             +\kappa_0(e_\Delta+e_{AG})\}.
 \label{eq:common-information-numerical-bound}
\end{equation}
Indeed, replacing both tables by their exact scalar functions at the
actual arguments costs their summed intrinsic errors. Exact decoded
addition gives \eqref{eq:common-information-numerical-endpoints}.
The effective gap need not be nonnegative, but the path bound in
Theorem~\ref{thm:common-information-stability} remains valid and the
final signed projection fixes the true correction. In the retained
graphs $B_A-B_G$ is dyadic and exactly encoded; retaining $e_{AG}$ in
the rule makes the required premise explicit. There is no sampled
inner Gaussian or strike law in this implementation, so adding $D_N$
or $\beta_S$ would charge an approximation it does not perform.
Normalisers other than one require separately certified scaling
coefficient and floor allowances and are not covered by this specialised
numerical formula.

\subsection{Matched finite shared-threshold implementation}

For a fixed ideal uniform strike midpoint $s$, put $a(s)=\log(1+s/100)$.
The exact graph computes a lower encoded endpoint, its upper endpoint
and one shared threshold as
\begin{equation}
 \widehat\ell=\widehat M_G+\widehat C,\qquad
 \widehat u=\widehat\ell+\widehat\Delta,\qquad
 \widehat z=\widehat a(s)+\widehat O
                       -\mathsf S_\lambda(\widetilde N+L).
 \label{eq:common-information-threshold-endpoints}
\end{equation}
Its predicate is $\widehat\ell<\widehat z\leq\widehat u$ and the gap
word is nonnegative. The represented strike error
$|\widehat a(s)-a(s)|\leq e_a$ includes its input-ratio arithmetic
and intrinsic logarithm error; it does not include normal quantile
rounding. The same $C,O,E_s$ as above identify the conceptual threshold
$a(s)-\lambda\widetilde N-\zeta(\widetilde N)$ relative to a common
represented lower endpoint. The represented strike and common endpoint
errors give a fixed-strike common shift bounded by
\begin{equation}
 E_{m,a}=e_m+e_C+e_O+e_a.
 \label{eq:common-information-threshold-common-error}
\end{equation}
They may depend arbitrarily on the fixed outer and strike labels.
They are independent of the inner normal label, which is the required
conditioning condition. With the retained strike midpoint bias $\beta_S$
in probability units, the implemented conditional event probability
obeys
\begin{equation}
 |p_{\rm thr}(y)-R_1(y)|\leq\eta_{\rm thr}:=
 \min\{1,\beta_S+2D_N+2\kappa_0E_s
                     +C(\Delta_*)E_{m,a}+\kappa_0e_\Delta\}.
 \label{eq:common-information-threshold-bound}
\end{equation}

For completeness, at a fixed strike and outer coordinate the bounded
$\zeta$ moves each of the two strict/non-strict Gaussian boundaries by
at most $E_s$. The represented event is contained in the interval
obtained by expanding the unperturbed interval by $E_s$ at each end.
It contains the interval obtained by shrinking by $E_s$ at each end,
whenever that interval is nonempty. The probability of either fixed
interval under the finite normal law differs from the continuous normal
probability by at most $2D_N$: an interval probability is a difference
of two right-CDF or left-CDF values. The continuous density is bounded
by $\kappa_0$, so expansion or shrinkage changes probability by at most
$2\kappa_0E_s$. If the shrunken interval is empty, its original width is
at most $2E_s$ and its entire continuous probability is bounded by
$2\kappa_0E_s$. Thus the finite event transfer costs $2D_N+2\kappa_0E_s$ once.
After this transfer, each fixed-strike continuous interval is the
Gaussian correction with threshold measure $\delta_{a(s)}$, which has
mass one. The common shift and gap allowances follow from
Theorem~\ref{thm:common-information-stability}. Averaging those uniform
allowances over the equal-weight strike labels preserves them; replacing
that strike law by the ideal uniform law charges the separately retained
$\beta_S$. This proves \eqref{eq:common-information-threshold-bound}.
Finite atoms are retained, and no second inner Gaussian-box mass is
multiplied into the decoder.

The label-dependent $\zeta$ cannot be charged merely with
$C(\Delta_*)E_s$ by invoking Gaussian convolution: it is correlated with
the normal label and the conditioning premise of that invocation fails.
The preceding sandwich proves the required allowance without such an
independence assumption.

\subsection{Gap-sensitive finite-law allowances and their limits}

The exact capped exponential difference is unimodal. Its total variation
in the normal coordinate is
\begin{equation}
 \operatorname{TV}\{f(m+\Delta+\lambda\,\cdot)-f(m+\lambda\,\cdot)\}
 =2v(\Delta),\qquad
 v(\Delta)=\min\{1,101(1-e^{-\Delta})\}.
 \label{eq:common-information-gap-variation}
\end{equation}
The retained earlier piecewise derivation proves this exact expression.
Accordingly, the raw finite-law transfer can use
$2D_Nv(\widehat\Delta)$; a true gap bound alone must first be enlarged
by its represented gap error. If $U=e^\Delta-1\geq0$, then
$v(\Delta)=\min\{1,101U/(1+U)\}$. This is a nondecreasing concave
function of $U$. Given an independently certified finite outer moment
$\mathbb E U\leq\overline U$, Jensen's inequality and the Lipschitz
constant $101$ of $v$ yield
\begin{equation}
 \mathbb E v(\widehat\Delta)
 \leq\min\left\{1,
       \min\left\{1,\frac{101\overline U}{1+\overline U}\right\}
                   +101e_\Delta\right\}.
 \label{eq:common-information-gap-mean}
\end{equation}
This is an averaged outer-law certificate, not a uniform conditional
certificate or an automatic two-sided prepared-probability prior.
For the four retained winning bases the certified first moments are
approximately $0.10888$, $0.02563$, $0.10684$ and $0.10666$.
The inner minimum in \eqref{eq:common-information-gap-mean} equals one
in all four cases. This moment-only refinement therefore provides no
positive cost consequence in the retained study. A tighter polynomial
moment majorant would require its own proved domination, optimisation
record and complete prior/budget reconstruction.

All three prepared methods can use their corresponding common-information
formula with the same $\nu_0,\nu_1$, true gap premise and financial
information. For the implemented basket, $\Delta=\log(1+U)\leq U\leq U_*$
provides a price-independent rational choice $\Delta_*=U_*$. A smaller
certified gap bound may be used only when retained with its derivation.
Any smaller probability allowance must be propagated through the
unchanged analytic coefficient, external outer-law mass, continuous
financial transfer, finite mean prior, estimation accuracy, confidence
and complete clean circuit cost. A certificate improvement changes
neither the ideal financial price nor the generic amplitude-estimation
query law.

\section{Representation dominance for reflection-only sampling}
\label{sec:representation-sampler-dominance}

The following lemma separates an implementation cost statement from an
estimator comparison. Let $A_i$ be a completely clean financial
preparation with its adjoint, and let $S_{0,i}$ reflect the complete
invariant input-zero subspace. The good reflection is a single Clifford
$Z$ on the already computed flag. Write $a_i=T(A_i)$ and
$s_i=T(S_{0,i})$, counting $T$ and $T^\dagger$ together. The complete
adjoint costs $a_i$. Measurement, classical conjunction of the observed
input-zero bits and conditional Pauli resets require no $T$ gates under
the declared logical architecture; their non-$T$ counts and allocated
wire counts remain separately charged. Scratch may be excluded from
the zero reflection only after its cleanliness invariant has been
proved. This lemma assumes those preparations and reflections, and
does not supply a missing coherent state-access or cleanliness proof.
The declared charge adds the complete retained preparation, adjoint and
reflection templates. It applies no cancellation or optimisation across
their boundaries. A different composition or cross-template optimisation
must be recosted before applying the theorem; the identities are not a
lower bound on every circuit realisation of the same polynomial signal.

\begin{lemma}[All-degree reflection-sampler dominance]
\label{thm:representation-sampler-dominance}
For every integer $d\geq1$, the retained parity-specific reflection-only
sampler has complete per-shot logical $T$ count
\begin{equation}
 t_i(d)=
 \begin{cases}
   da_i+(d-1)s_i/2,&d\text{ odd},\\
   da_i+(d/2-1)s_i,&d\text{ even}.
 \end{cases}
 \label{eq:representation-parity-cost}
\end{equation}
Candidate $c$ has $t_c(d)\leq t_b(d)$ for every positive integer $d$
if and only if
\begin{equation}
 a_c\leq a_b,\qquad 2a_c+s_c\leq2a_b+s_b.
 \label{eq:representation-all-degree-condition}
\end{equation}
If $a_b>0$ and $2a_b+s_b>0$, then every degree obeys the uniform ratio
bound
\begin{equation}
 \frac{t_c(d)}{t_b(d)}\leq
 \max\left\{\frac{a_c}{a_b},
             \frac{2a_c+s_c}{2a_b+s_b}\right\}.
 \label{eq:representation-all-degree-ratio}
\end{equation}
Both endpoint ratios are necessary for a uniform bound of this form:
the first is attained at $d=1,2$ and the second is approached as
$d\longrightarrow\infty$. Strict superiority for every degree holds
if and only if $a_c<a_b$ and $2a_c+s_c\leq2a_b+s_b$.
\end{lemma}
\begin{proof}
Put $Q_i=-A_iS_{0,i}A_i^\dagger S_g$. The unconditional global minus
can be omitted because it changes a terminal shot only by a global
phase. It cannot be omitted from a controlled-QPE comparison. For
$d=2k+1$, the chronological sampler prepares $A_i$ and then applies
$Q_i$ exactly $k$ times. It uses $k+1$ preparations, $k$ adjoints,
$k$ zero reflections and $k$ Clifford good reflections. For $d=2k$
with $k\geq1$, it prepares $A_i$, applies $Q_i^{k-1}$, applies $S_g$
and then $A_i^\dagger$, observing the invariant input-zero event. It
uses $k$ preparations, $k$ adjoints, $k-1$ zero reflections and $k$
Clifford good reflections. Complete measurement and reset are appended
to both constructions. Their declared $T$ charge is zero, so these
counts prove \eqref{eq:representation-parity-cost}.

Let $q_i=2a_i+s_i$, $\delta a=a_c-a_b$ and $\delta q=q_c-q_b$.
The odd costs are $a_i+kq_i$ for $k\geq0$ and the even costs are
$2a_i+(k-1)q_i$ for $k\geq1$. If $\delta a\leq0$ and
$\delta q\leq0$, subtracting these formulae gives nonpositive
differences in both parities. Conversely, degree one requires
$\delta a\leq0$. If $\delta q>0$, the odd difference
$\delta a+k\delta q$ is positive at a sufficiently large integer
$k$, contradicting all-degree dominance. This proves necessity and
sufficiency. At $d=1$ strict dominance requires $\delta a<0$; together
with $\delta q\leq0$ this is also sufficient in both parities.

For odd degrees the ratio of costs is
$(a_c+kq_c)/(a_b+kq_b)$. It is a convex combination of $a_c/a_b$
and $q_c/q_b$, with nonnegative weights
$a_b/(a_b+kq_b)$ and $kq_b/(a_b+kq_b)$. For even degrees replace
$a_i$ by $2a_i$ and $k$ by $k-1$; the two endpoint ratios are
unchanged. A convex combination is bounded by the larger endpoint.
Degrees one and two attain $a_c/a_b$. Along either parity the ratio
tends to $q_c/q_b$ as its degree tends to infinity, proving the final
uniform sharpness statement.
\end{proof}

\begin{corollary}[Same-history adaptive sampler comparison]
\label{cor:representation-same-history}
Suppose two clean representations have the same implemented good
probability $p$, and an estimator uses the same initial prior, tolerance,
confidence allocation and deterministic history-dependent degree/batch
policy for both. Under ideal logical operation their complete reflection
shots can be coupled to have identical observed Bernoulli outcomes,
chosen degrees, batch sizes and stopping decisions. If
\eqref{eq:representation-all-degree-condition} holds, every coupled
history has no larger candidate reflection-shot $T$ cost. If
\eqref{eq:representation-all-degree-ratio} has right side $r$, then
the total candidate reflection-shot cost on every history is at most
$r$ times the corresponding rival reflection-shot cost. Consequently
the same comparison holds for expectations whenever they exist.
\end{corollary}
\begin{proof}
The two-dimensional reflection algebra gives the same
Chebyshev Bernoulli parameter $q_d(p)$ for both clean representations
at each degree. This includes $p=0,1$ by direct endpoint evaluation,
without normalising a zero-probability basis component. At each shot,
use a fresh common uniform random variable $V$ to define both observed
bits as $\mathbf1_{\{V\leq q_d(p)\}}$. Inductively, identical preceding
histories imply the same next degree and batch size, which imply
identical next observation under this coupling. The deterministic
policy then makes the same next decision. Apply the pointwise
per-shot inequality to each degree in the common history and sum its
nonnegative costs. Taking expectation preserves the inequality.
For unbounded stopping histories, apply the same argument to finite
prefixes and then use monotone convergence. Any additional fallback
or non-reflection primitive requires its own cost comparison; it is not
included by this corollary merely because it belongs to the estimator.
\end{proof}

The continuous ideal numerical, raw-payoff and shared-threshold
constructions target the same financial mean and hence the same ideal
Grover angle after a matched monetary normaliser. Their separately
rounded finite implementations need not have identical $p$. Their
deterministic allowances, probability priors, admissible degrees and
stopping histories can therefore differ. The all-degree cost theorem
compares the actual complete preparation resources at matched degrees;
the same-history corollary is conditional on matched implemented
probability and controller information. Neither theorem alone proves
smaller realised adaptive costs, a new amplitude-estimation query law,
dominance over non-reflection estimators or physical quantum advantage.
The generic reflection algebra is established. The project-specific
positive consequence is an exact, fully charged representation
comparison if the retained preparations satisfy its two affine
inequalities; mathematical priority remains a separate literature
question.

% Supporting fragment: established confidence tools at a matched financial interface.
\subsection{Matched financial guarantees and paired diagnostic signals}
\label{supp:matched-estimator-transfer}

Let $V$ be the single stipulated derivative value. For alternative
$i$, let its certified clean logical preparation $A_i$ have actual
success probability $p_i$ and deterministic prior
$[l_i,u_i]\subseteq[0,1]$. Suppose its certified monetary decoding
$D_i(p)$ satisfies
\begin{equation}
 |V-D_i(p_i)|\leq E_i,
 \qquad |D_i(p)-D_i(q)|\leq B_i|p-q|.
 \label{eq:matched-financial-interface}
\end{equation}
The certificate includes every deterministic representation, arithmetic,
geometric-control and external-factor contribution charged by that
alternative. A common target radius $\varepsilon_V$ is met whenever
\begin{equation}
 E_i+B_i\varepsilon_i\leq\varepsilon_V.
 \label{eq:matched-monetary-budget}
\end{equation}
Thus monetary matching does not require equal probability tolerances:
different deterministic errors or normalisers can require different
$\varepsilon_i$. Replacing them by a common larger tolerance would need
a new monetary certificate.

For each alternative the existing bounded controller uses at most
$L=12$ fresh batches of $n=100$ shots. The pilot spends
$\alpha/(2L)$ at every reached look, including repeated degrees; the
fresh original-prior fallback receives $\alpha/2$. Conditional on that
alternative's preceding history, the next selected degree is fixed and
its new shots are independent Bernoulli measurements of the exact
logical signal. The fixed-sample binomial interval therefore covers
conditionally with probability at least $1-\alpha/(2L)$.
The union bound over the at most $L$ reached looks gives pilot failure
at most $\alpha/2$. Adding the fallback's failure allowance proves
\begin{equation}
 \Pr\{|\widehat V_i-V|>\varepsilon_V\}\leq\alpha,
 \label{eq:matched-alternative-coverage}
\end{equation}
under the ideal logical-shot premises and the separately certified
deterministic interface. An empty intersection is an unsuccessful pilot,
which invokes the charged fallback using the original prior. The pilot
does not supply an uncertified tightened prior to that fallback.

Let $T_{F,i}$ be the complete fresh half-$\alpha$ canonical fallback
cost and $C_i=\lfloor T_{F,i}/10\rfloor$ the pilot cap. The controller
checks a full batch charge before requesting it. With $S_i$ denoting
successful pilot termination and $T_{P,i}$ its complete pilot charge,
every possible history satisfies
\begin{equation}
 T_i=T_{P,i}+T_{F,i}\mathbf1_{S_i^c},
 \qquad 0\leq T_{P,i}\leq C_i,
 \qquad T_i\leq C_i+T_{F,i}.
 \label{eq:matched-history-cap}
\end{equation}
Every preparation, adjoint, reflection, observation and reset contributes
to these charges. The half-$\alpha$ fallback can require more repetitions
than the full-$\alpha$ canonical rival $T_{0,i}$. Consequently
$C_i+T_{F,i}$ is not generally bounded by $1.1T_{0,i}$.
Taking expectations in \eqref{eq:matched-history-cap} requires no
independence between success and pilot cost and yields the exact criterion
\begin{equation}
 \mathbb E T_i<T_{0,i}
 \quad\Longleftrightarrow\quad
 \Pr(S_i)>1-\frac{T_{0,i}}{T_{F,i}}
                 +\frac{\mathbb E T_{P,i}}{T_{F,i}}.
 \label{eq:matched-break-even}
\end{equation}
A finite equal-weight mesh supplies a descriptive version of this
identity. It does not identify either expectation or success probability
at an unknown actual financial flag.

For the paired diagnostic, choose a stipulated $p_*$ in a nonempty
intersection of all alternative priors and their common finite ideal
mean interval. Let $U_{j,s}$ be independent uniform variables for
look $j$ and shot $s$, reused across alternatives. If alternative $i$
selects degree $d_{i,j}$ from its own previous outcomes, define
\begin{equation}
 X_{i,j,s}=\mathbf1_{\{U_{j,s}<q_{d_{i,j}}(p_*)\}},
 \qquad
 q_d(p)=T_d(\sqrt p)^2.
 \label{eq:matched-paired-uniforms}
\end{equation}
For a fixed alternative, $U_{j,s}$ is independent of that alternative's
earlier looks. Its conditional fresh-shot law is therefore exactly the
law required by the controller. The alternatives need not be independent
for their marginal confidence bounds or the union bound to hold.
This reasoning does not condition a method on another method's entire
future history. At $\alpha=1/3600$, all twelve alternative-price
events have union failure at most $1/300$; selecting one independently
certified alternative for each of four cases gives $1/900$.

The finite replay uses pseudorandom64-bit words and rounded signal
thresholds solely to reproduce descriptive histories. Its word coupling
removes method-relative mesh and seed confounding. It supplies no
theorem about pseudorandom or physical coverage. In particular,
\eqref{eq:matched-paired-uniforms} stipulates a common scalar signal;
it does not assert $p_i=p_j$ for differently rounded financial
preparations. Their actual flags remain governed by their own
representation certificates. Shared ideal financial means and distinct
implemented flag probabilities are compatible statements.

These arguments specialise established fixed-binomial coverage,
conditional confidence spending, union bounds and linear expectation to
the project interface. I do not claim their generic theoretical novelty.

\section{Sharp Gaussian component moduli}
\label{sec:sharp-common-information}

The constants below refine the existing path certificate without using a
financial price as input. They are exact extrema for the stated class of
positive mass-one threshold measures. Their mathematical sharpness does
not establish priority in the literature. In particular, the two
component extrema need not be attained by the same perturbation; the
sum in the path certificate is not asserted to be jointly sharp.

Write $\phi(x)=(2\pi)^{-1/2}\exp(-x^2/2)$ and let $\Phi$ be its
distribution function. For $t\geq0$, define
\begin{equation}
 k(t)=\sup_{x\in\mathbb R}|\phi(x+t)-\phi(x)|,
 \qquad \mathfrak a(t)=2\Phi(t/2)-1.
 \label{eq:sharp-k-A}
\end{equation}
For $t,e\geq0$, put $r_t(x)=\Phi(x+t/2)-\Phi(x-t/2)$ and
\begin{equation}
 \Omega(t,e)=\sup_{x\in\mathbb R,\,|h|\leq e}
                    |r_t(x+h)-r_t(x)|.
 \label{eq:sharp-Omega-definition}
\end{equation}
Both $k(0)$ and $\Omega(t,0)=\Omega(0,e)$ are zero.

\begin{theorem}[Exact derivative coefficient]
\label{thm:sharp-common-derivative}
For $t>0$, there is exactly one $u_t>\max\{1,t/2\}$ satisfying
\begin{equation}
 u_t\tanh(tu_t/2)=t/2,
 \qquad
 k(t)=\phi(u_t-t/2)-\phi(u_t+t/2).
 \label{eq:sharp-k-root}
\end{equation}
It lies below $1+t/2$. The function $k$ is strictly increasing on
$(0,\infty)$ and
\begin{equation}
 0<k(t)<\min\{\phi(0),t\phi(1)\},\qquad
 \lim_{t\downarrow0}\frac{k(t)}{t}=\phi(1),\qquad
 \lim_{t\to\infty}k(t)=\phi(0).
 \label{eq:sharp-k-properties}
\end{equation}
For every positive Borel measure $\mu$ of mass one and every $\lambda>0$,
the smallest uniform coefficient, simultaneously over $m$ and $\mu$, in
$|\partial_m R_{\lambda,\mu}(m,d)|\leq\kappa(d,\lambda)$ is
\begin{equation}
 \kappa(d,\lambda)=\lambda^{-1}k(d/\lambda).
 \label{eq:sharp-common-coefficient}
\end{equation}
An atomic measure attains this coefficient.
\end{theorem}
\begin{proof}
For $a=t/2$ and $u>0$, Gaussian symmetry gives the positive function
$q_t(u)=\phi(u-a)-\phi(u+a)$; the absolute density difference reduces
to this function by translation and reflection. Its logarithmic
derivative is
\begin{equation}
 \frac{q_t'(u)}{q_t(u)}=-u+a\coth(au).
 \label{eq:sharp-k-log-derivative}
\end{equation}
The function $u\mapsto u\tanh(au)$ is continuous, strictly increasing
from zero to infinity. At $u=a$ its value is strictly below $a$, and at
$u=1$ its value is strictly below $a$, since $\tanh a<a$. Thus its
unique level-$a$ point exceeds both. The inequality
$\tanh x\geq x/(1+x)$ for $x\geq0$ follows from
$\exp(2x)\geq1+2x$. At $u=a+1$ it implies
\begin{equation}
 u\tanh(au)\geq
 \frac{a(a+1)^2}{1+a(a+1)}>a.
 \label{eq:sharp-k-upper-bracket}
\end{equation}
Hence $q_t$ increases up to $u_t$ and decreases thereafter, proving
the formula and bracket. At the maximiser $u_t>a$, differentiation
with respect to $t$ gives
\begin{equation}
 \partial_t q_t(u_t)=\tfrac12\{(u_t-a)\phi(u_t-a)
                              +(u_t+a)\phi(u_t+a)\}>0.
 \label{eq:sharp-k-monotonicity}
\end{equation}
The implicit function theorem applies because
$\partial_u[u\tanh(au)]>0$, so this is the derivative of $k$.
Both Gaussian densities are positive, giving $k(t)<\phi(0)$.
The fundamental theorem of calculus gives
$q_t(u)\leq t\sup|\phi'|=t\phi(1)$, strictly for a positive interval,
since the continuous function $|\phi'|$ attains its maximum only at
two isolated points. The bracket and root equation imply $u_t\to1$
as $t\downarrow0$, establishing the small-$t$ limit. Evaluating the
density difference at $x=0$ gives the lower bound
$\phi(0)-\phi(t)$, which proves the large-$t$ limit. Finally,
\begin{equation}
 \partial_m R_{\lambda,\mu}(m,d)=\lambda^{-1}
   \int\{\phi((m+d-a)/\lambda)-\phi((m-a)/\lambda)\}\,\mu(da).
 \label{eq:sharp-common-derivative-integral}
\end{equation}
Positivity and total mass one prove the coefficient bound. Choosing
$\mu=\delta_a$ and $m-a=\lambda(u_t-t/2)$ attains its absolute value.
\end{proof}

\begin{theorem}[Exact finite common-shift modulus]
\label{thm:sharp-finite-common-shift}
For $t,e>0$, put $b=(t+e)/2$ and $c=|t-e|/2$. There is exactly one
$u_{t,e}>\max\{t,e\}/2$ satisfying
\begin{equation}
 \log\cosh(bu_{t,e})-\log\cosh(cu_{t,e})=te/2.
 \label{eq:sharp-Omega-root}
\end{equation}
It is smaller than $\max\{t,e\}/2+1/\min\{t,e\}$, and
\begin{equation}
 \Omega(t,e)=r_t(u_{t,e}-e/2)-r_t(u_{t,e}+e/2).
 \label{eq:sharp-Omega-evaluation}
\end{equation}
The function $\Omega$ is symmetric in its arguments and strictly
increasing in each positive argument. Moreover,
\begin{equation}
 0<\Omega(t,e)<\min\{k(t)e,k(e)t,\mathfrak a(t),\mathfrak a(e)\}.
 \label{eq:sharp-Omega-comparison}
\end{equation}
For every positive mass-one $\mu$,
\begin{equation}
 |R_{\lambda,\mu}(m+h,d)-R_{\lambda,\mu}(m,d)|
 \leq\Omega(d/\lambda,E/\lambda),\qquad |h|\leq E, d\geq0.
 \label{eq:sharp-Omega-mixture}
\end{equation}
The right-hand side is the smallest uniform bound over this class,
and is attained by a point mass and a shift of magnitude $E$.
\end{theorem}
\begin{proof}
The function $r_t$ is even, strictly decreasing on $(0,\infty)$ and
positive everywhere for $t>0$. For a shift of magnitude $e$, reflection
reduces the absolute difference to
$D_{t,e}(u)=r_t(u-e/2)-r_t(u+e/2)$ for $u\geq0$. This quantity is
positive for $u>0$. Expanding the four distribution functions shows
$D_{t,e}=D_{e,t}$. Differentiation gives
\begin{equation}
 D_{t,e}'(u)=2\phi(u)\{
 e^{-c^2/2}\cosh(cu)-e^{-b^2/2}\cosh(bu)\}.
 \label{eq:sharp-Omega-derivative}
\end{equation}
Since $b>c\geq0$, the function
$H(u)=\log\cosh(bu)-\log\cosh(cu)$ is strictly increasing on
$(0,\infty)$: its derivative is
$b\tanh(bu)-c\tanh(cu)>0$, because $v\mapsto v\tanh(vu)$ is
strictly increasing for $v>0$. Also $H(0)=0$ and $H(u)\to\infty$.
Since $(b^2-c^2)/2=te/2$, the unique root in
\eqref{eq:sharp-Omega-root} gives the unique maximum of $D_{t,e}$.
At $u=e/2$, the derivative is $r_t'(0)-r_t'(e)>0$; symmetry also
gives positivity at $u=t/2$. This proves the lower bracket.
The inequalities $\cosh(bu)\geq e^{bu}/2$ and
$\cosh(cu)\leq e^{cu}$ imply
$H(u)\geq(b-c)u-\log2$. At the stated upper bracket this exceeds
$te/2$, since $b-c=\min\{t,e\}$ and $1>\log2$.

The maximum depends smoothly on $e>0$ by the implicit function
theorem. At its maximiser, $D_{t,e}'(u)=0$, so
$r_t'(u-e/2)=r_t'(u+e/2)$. Consequently
\begin{equation}
 \partial_e\Omega(t,e)
   =-r_t'(u_{t,e}-e/2)>0.
 \label{eq:sharp-Omega-monotonicity}
\end{equation}
The strict inequality uses $u_{t,e}>e/2$. It follows that the largest
modulus over all shifts of magnitude at most $e$ is attained at $e$.
Symmetry gives monotonicity in $t$ as well. Integrating $r_t'$ over
the interval of length $e$ gives $\Omega(t,e)<k(t)e$, strictly because
the continuous derivative cannot be constant at its extremal value on
a positive interval. Interchanging $t,e$ gives $\Omega(t,e)<k(e)t$.
The largest value of $r_t$ is $r_t(0)=\mathfrak a(t)$, whereas the smaller
of the two values in the attained difference remains positive; hence
$\Omega(t,e)<\mathfrak a(t)$ and symmetry gives $\Omega(t,e)<\mathfrak a(e)$.
For a fixed threshold $a$, the correction is
$r_{d/\lambda}((m-a+d/2)/\lambda)$. Integrating its common-shift
bound against $\mu$ proves \eqref{eq:sharp-Omega-mixture}.
An atomic measure with the lower centred argument
$u_{t,e}-e/2$ and shift $E$ attains the absolute difference.
\end{proof}

\begin{corollary}[Finite path certificate]
\label{cor:sharp-finite-path}
Under the premises of Theorem~\ref{thm:common-information-stability},
\begin{equation}
 |R_{\lambda,\mu}(\widehat m,\widehat\Delta)
                        -R_{\lambda,\mu}(m,\Delta)|
 \leq\Omega(\Delta/\lambda,E_m/\lambda)
                       +\mathfrak a(e_\Delta/\lambda).
 \label{eq:sharp-finite-path}
\end{equation}
The bound remains valid when the represented gap is negative and after
projection that fixes the true correction. A known upper true gap may
replace $\Delta$. If both gaps are nonnegative, the first argument of
$\Omega$ may instead use $\min\{\Delta,\widehat\Delta\}$.
\end{corollary}
\begin{proof}
Use the same two-segment path as the earlier theorem. The common-shift
segment receives \eqref{eq:sharp-Omega-mixture}. For the gap segment,
the integrand is the Gaussian probability of an interval of length
$|\widehat\Delta-\Delta|/\lambda$. The probability of an interval
of length $v$ is maximised by the centred interval, as follows by
differentiating $\Phi(x+v)-\Phi(x)$ and using Gaussian symmetry.
Its maximum is $\mathfrak a(v)$. This argument also covers a negative represented
gap. Positivity and mass one preserve the bound under integration.
The monotonicity of $\mathfrak a$ and $\Omega$ gives the stated replacements;
reversing the path gives the smaller represented-gap alternative.
Projection is $1$-Lipschitz and fixes the target.
\end{proof}

\subsection{Certified evaluation without a financial price input}

All inputs to $k$, $\mathfrak a$ and $\Omega$ are the smoothing scale and
independently certified error/gap bounds. The root equations have unique
roots, explicit brackets admitting rational outward enclosures and
strictly monotone left-hand sides. Directed interval evaluation and rational bisection
therefore enclose the roots and the constants to arbitrary precision.
When an interval evaluation cannot decide the sign, increasing arithmetic
precision decides it unless the rational trial point is the exact root;
that case is covered as follows. For a retained root bracket $[l,r]$,
a trial point $c\in[l,r]$, a certified residual $|F(c)|\leq\varepsilon$
and a certified derivative lower bound
$\inf_{[l,r]}F'\geq a>0$, the mean value theorem gives the new enclosure
\begin{equation}
 u_*\in[l,r]\cap[c-\varepsilon/a,c+\varepsilon/a].
 \label{eq:sharp-root-residual-enclosure}
\end{equation}
The analytic derivative is strictly positive on the stated lower brackets;
a sufficiently fine interval subdivision certifies a positive lower
bound on the compact retained bracket. Evaluating the expressions on the entire retained root interval
gives an upper bound, rather than treating a rounded root as exact.
Outward rational rounding may subsequently preserve the upper bound.
Neither a realised derivative price nor a simulated price enters this
procedure. The constants are applicable to every encoding satisfying the
same conditioning premise. They do not justify treating normal-label
dependent arithmetic perturbations as independent common shifts.

\subsection{The actual capped-payoff threshold measure}

For $f(x)=H^{-1}\min\{K(e^x-1)_+,H\}$, write
$\alpha=K/H$ and $\ell=\log(1+H/K)$. Its threshold measure is
\begin{equation}
 \mu_{\rm cap}(da)=\alpha e^a\mathbf1_{(0,\ell)}(a)\,da,
 \qquad\int\mu_{\rm cap}(da)=\alpha(e^\ell-1)=1.
 \label{eq:sharp-cap-measure}
\end{equation}
Its smoothed derivative, which is a density rather than a basket price,
has the explicit expression
\begin{equation}
 j_\lambda(m)=P_{\lambda,\mu_{\rm cap}}'(m)
 =\alpha e^{m+\lambda^2/2}
 \left\{\Phi\left(\frac{\ell-m-\lambda^2}{\lambda}\right)
              -\Phi\left(\frac{-m-\lambda^2}{\lambda}\right)\right\}.
 \label{eq:sharp-cap-density}
\end{equation}
Define the model-specific coefficient
\begin{equation}
 \kappa_{\rm cap}(d,\lambda)
  =\sup_{m\in\mathbb R}|j_\lambda(m+d)-j_\lambda(m)|.
 \label{eq:sharp-cap-kappa}
\end{equation}
For every $d>0$ and $\lambda>0$,
\begin{equation}
 0<\kappa_{\rm cap}(d,\lambda)
   <\kappa(d,\lambda),\qquad
 \kappa_{\rm cap}(d,\lambda)\leq(K+H)/H.
 \label{eq:sharp-cap-strict-improvement}
\end{equation}
To derive \eqref{eq:sharp-cap-density}, integrate
$\alpha e^a\lambda^{-1}\phi((m-a)/\lambda)$ over $(0,\ell)$
and complete the square in $a$. The result is
\eqref{eq:sharp-cap-density}. The function $j_\lambda$ is continuous,
strictly positive and tends to zero at both tails. Its translated
difference is not identically zero for $d>0$: otherwise
$j_\lambda(m+d)=j_\lambda(m)$ for every $m$, making this positive
function periodic, contrary to its zero tail limit. Thus its translated
difference attains its nonzero absolute maximum at a finite point.
For each fixed $m$, the Gaussian translated-density difference attains
its extremal absolute value at only two isolated threshold positions.
The positive absolutely continuous measure \eqref{eq:sharp-cap-measure}
gives those positions zero mass. Strictness of the integral triangle
bound at the attained maximiser proves the strict inequality in
\eqref{eq:sharp-cap-strict-improvement}. Finally,
$j_\lambda=\mathbb E f'(m+\lambda N)$ and
$0\leq f'\leq(K+H)/H$, so the difference of two such values has
absolute value at most $(K+H)/H$.

This coefficient is certifiably computable. For any rational $L>0$, all
$m\notin[-d-L\lambda,\ell+L\lambda]$ obey the tail bound
\begin{equation}
 |j_\lambda(m+d)-j_\lambda(m)|\leq\phi(L)/\lambda.
 \label{eq:sharp-cap-tail}
\end{equation}
Indeed, both Gaussian arguments then remain on the same side of zero
with magnitude at least $L$ for every threshold in $[0,\ell]$.
Their nonnegative densities each lie below $\phi(L)$, and their
absolute difference lies below the same bound. On the compact interval,
the integrand derivative gives
\begin{equation}
 |\partial_m\{j_\lambda(m+d)-j_\lambda(m)\}|
                         \leq2\phi(1)/\lambda^2.
 \label{eq:sharp-cap-mesh-Lipschitz}
\end{equation}
A mesh of maximum spacing $\delta$, with directed upper bounds $U_i$
for the absolute differences at its nodes and both endpoints included,
therefore yields the global certificate
\begin{equation}
 \kappa_{\rm cap}(d,\lambda)
 \leq\max\left\{\phi(L)/\lambda,
                       \max_i U_i+\phi(1)\delta/\lambda^2\right\}.
 \label{eq:sharp-cap-mesh-certificate}
\end{equation}
Letting $L$ increase, $\delta$ decrease and evaluation precision increase
converges to the exact supremum. A single coefficient at $d_*$ cannot be
used for every $0\leq d\leq d_*$ without a separate monotonicity argument.
A directly certified alternative is
\begin{equation}
 \overline\kappa_{\rm cap}(d_*,\lambda)
 =\sup_{0\leq v\leq d_*}\sup_m
                  |j_\lambda(m+v)-j_\lambda(m)|
 <\kappa(d_*,\lambda)\qquad(d_*>0).
 \label{eq:sharp-cap-uniform-gap-envelope}
\end{equation}
For its strict inequality, the joint function is continuous and tends
uniformly to zero in the two tails for $v\in[0,d_*]$. Its positive
maximum is attained on a compact set. At an attained pair with $v>0$,
the preceding strict cap inequality gives a value smaller than
$\kappa(v,\lambda)\leq\kappa(d_*,\lambda)$; at $v=0$ the value is zero.
To compute this envelope, retain the compact $m$ interval
$[-d_*-L\lambda,\ell+L\lambda]$ and a rectangular mesh with maximum
spacings $\delta_m,\delta_v$. Since the derivatives in the two
coordinates have absolute bounds $2\phi(1)/\lambda^2$ and
$\phi(1)/\lambda^2$, respectively, the maximum node upper bound plus
$\phi(1)(\delta_m+\delta_v/2)/\lambda^2$, combined with the tail
bound \eqref{eq:sharp-cap-tail}, is a global upper certificate.
Any use of this model-specific refinement
must preserve the represented conditioning information. In particular,
the shared-threshold implementation may have strike-label dependent
threshold errors. It cannot apply the integrated cap coefficient to
those errors. The strike-independent common endpoint shift may use the
cap coefficient; the remaining strike-dependent shift must retain a
uniform fixed-threshold allowance. Raw-payoff and numerical encodings
may use the cap coefficient only after their common endpoint errors have
been separated from scalar-map errors and inner-label dependent
remainders as in the existing guarded certificates.

\subsection{Finite remainders that depend on the inner label}

These allowances remain separate from common-information smoothing.
For any Borel set $B$ of Lebesgue measure at most $w\geq0$,
\begin{equation}
 \int_B\phi(x)\,dx\leq \mathfrak a(w).
 \label{eq:sharp-normal-set-concentration}
\end{equation}
Indeed, let $I=[-w/2,w/2]$. The density is at most $\phi(w/2)$ on
$B\setminus I$ and at least that value on $I\setminus B$. Thus
\begin{equation}
 \int_B\phi-\int_I\phi
 \leq\phi(w/2)\{|B\setminus I|-|I\setminus B|\}\leq0.
 \label{eq:sharp-normal-set-proof}
\end{equation}
The two boundary strips in a Gaussian interval expansion or shrinkage
by $E_s/\lambda$ have total length at most $2E_s/\lambda$.
If shrinkage empties the interval, its original length is at most this
same bound. The existing interval sandwich therefore sharpens the
label-dependent shared-threshold allowance to
\begin{equation}
 2D_N+\mathfrak a(2E_s/\lambda)
       \leq2D_N+2\phi(0)E_s/\lambda.
 \label{eq:sharp-threshold-inner-remainder}
\end{equation}
Only the continuous strip probability is bounded by concentration;
the finite law still uses its retained interval CDF discrepancy and
strict/left endpoint conventions. This argument imposes no independence
between the perturbation and the inner normal label.

The normalised cap has the exact finite payoff modulus
\begin{equation}
 \omega_f(h)=\sup_x|f(x+h)-f(x)|
       =\min\{1,(\alpha+1)(1-e^{-|h|})\}.
 \label{eq:sharp-cap-payoff-modulus}
\end{equation}
For $h\geq0$, the increment integrates $\alpha e^x$ over the intersection
of an interval of length $h$ with $(0,\ell)$. Since that density is
increasing inside its support, its largest integral uses the rightmost
subinterval of length $\min\{h,\ell\}$. This proves the formula;
negative shifts give the same modulus by translation. For $d\geq0$,
the difference $g_d(x)=f(x+d)-f(x)$ lies in $[0,v(d)]$, where
$v(d)=\min\{1,(\alpha+1)(1-e^{-d})\}$. If $h\geq0$, both increments
$f(x+d+h)-f(x+d)$ and $f(x+h)-f(x)$ lie in $[0,\omega_f(h)]$.
Their difference consequently has magnitude at most $\omega_f(h)$.
The same argument follows for $h<0$ by translation. Hence
\begin{equation}
 |g_d(x+\zeta)-g_d(x)|
 \leq\min\{v(d),\omega_f(E_s)\}\qquad(|\zeta|\leq E_s).
 \label{eq:sharp-raw-inner-remainder}
\end{equation}
This replaces the former $(\alpha+1)E_s$ allowance and is valid
pointwise, including arbitrary dependence on the inner label.

\section{Exact conjugation and complete charging of an unclean preparation}
\label{app:cancellation-aware-preparation}

The preparation-conjugated reflection and the parity-dependent sampling
law are established amplitude-amplification constructions
\citep{suzuki2020,rallfuller2023}. Cross-boundary cancellation is also
established: Stamatopoulos et al.\ optimise the complete pricing circuit
in Appendix B, and the author version of Carrera Vazquez and Woerner
uses a spin-echo identity to reduce repeated preparation costs in
Section IV, equations (19)--(25), and Figure 2.
\footnote{The published record is \emph{Physical Review Applied}, 15,
034027, DOI 10.1103/PhysRevApplied.15.034027. The detailed text actually
consulted is arXiv:2005.07711v1, dated 15 May 2020, with a body date of
19 May 2020. Its rotation-specific argument is not asserted to be
identical to the inaccessible final text.}
I derive the following identities for the exact finite financial
circuits and charge the same transformation to every financial rival.
They justify an implementation comparison, without a claim of priority
for conjugation or amplitude estimation.

\subsection{Registers, operator order and the transported invariant}

Let $\mathcal H=\mathcal H_C\otimes\mathcal H_W$ be the financial
Hilbert space. The core $C$ contains the original uniform input labels,
any original comparison label and the good flag $g$. The work register
$W$ contains arithmetic outputs, private evaluator work, constant-zero
endpoint wires and any clean predicate scratch. Define
\begin{equation}
 \mathcal M=\mathcal H_C\otimes\operatorname{span}\{|0_W\rangle\},
 \qquad |e\rangle=|0_C\rangle\otimes|0_W\rangle,
 \qquad E_C=|0_C\rangle\langle0_C|\otimes I_W.
 \label{eq:conjugation-clean-space}
\end{equation}
The desired logical zero reflection is the \emph{same} old core operator
\begin{equation}
 S_0=I-2E_C,
 \qquad S_g=Z_g=I-2\Pi_g,
 \qquad \Pi_g=|1_g\rangle\langle1_g|\otimes I_{\mathrm{other}}.
 \label{eq:conjugation-reflections}
\end{equation}
Let $\widetilde S_0$ be the exact full unitary implemented by the
retained ladder tape on its actual allocated wires. The tape may borrow
scratch wires from $W$, and its established contract is
\begin{equation}
 \widetilde S_0x=S_0x\quad(x\in\mathcal M).
 \label{eq:conjugation-physical-zero-reflection}
\end{equation}
Both versions retain this \emph{same actual tape}. On $\mathcal M$,
its action is the rank-one reflection about $|e\rangle$, and all
borrowed scratch is restored to zero. I make no assumption that its
action equals the logical extension $I-2E_C$ on arbitrary dirty work.
That extension would itself not be rank one on the whole Hilbert
space. The distinction is necessary when asserting a full-operator
conjugacy: its two sides must contain the same actual unitary, rather
than replacing an invariant-space implementation by an unproved
global extension.
When a controlled ladder requires an additional scratch wire, I include
it in $W$ and extend $A$, $B$, $U$ and the uncontrolled tape by the
identity on that wire. The same zero-work space is thereby available
for both constructions on a common conservative allocation.

Let $H$ denote the exact tensor product of the retained input
Hadamards, let $U$ compute the reversible arithmetic graph, and let
$P$ be the exact clean XOR of its final predicate into $g$. Operator
products act from right to left, so the retained chronological clean
preparation is $H,U,P,U^\dagger$. I require the following checked
finite-circuit properties:
\begin{equation}
 A=U^\dagger P U H,\qquad
 A\mathcal M\subseteq\mathcal M,\qquad
 A^\dagger\mathcal M\subseteq\mathcal M,\qquad
 [U,S_g]=0.
 \label{eq:conjugation-assumptions}
\end{equation}
The last property follows if no gate in $U$ acts on the flag. The
cleanliness properties require $P$ to preserve the arithmetic output
registers and return its own private scratch to zero, so that the
emitted $U^\dagger$ restores every graph work register. Those properties
must hold on the whole registered zero-work core subspace, including
both flag values, rather than only on the single prepared input.
Exact circuit reversibility is compatible with approximate financial
tables: a table and its exact emitted adjoint cancel as finite
operators, irrespective of their approximation error relative to a
continuous financial function.

Define the alternative preparation
\begin{equation}
 B=P U H=U A,\qquad
 |\psi\rangle=A|e\rangle,\qquad
 |\phi\rangle=B|e\rangle=U|\psi\rangle.
 \label{eq:conjugation-preparations}
\end{equation}
The state $|\phi\rangle$ need not have clean arithmetic work. It is
therefore incorrect to replace its transported invariant by
$\mathcal M$ or to infer a reduction in the reflection core or allocated
peak from the removal of its final cleanup.

\begin{theorem}[Exact signal-preserving preparation conjugation]
\label{thm:exact-preparation-conjugation}
Under \eqref{eq:conjugation-assumptions} and the actual-tape clean-work
contracts \eqref{eq:conjugation-physical-zero-reflection} and
\eqref{eq:conjugation-actual-controlled-contract}, retain the
conditional sign in
\begin{equation}
 Q_A=-A\widetilde S_0 A^\dagger S_g,\qquad
 Q_B=-B\widetilde S_0 B^\dagger S_g.
 \label{eq:conjugation-iterates}
\end{equation}
The global unitary identity is
\begin{equation}
 Q_B=U Q_A U^\dagger,\qquad
 Q_B^j B|e\rangle=UQ_A^j A|e\rangle\quad(j\geq0).
 \label{eq:conjugation-global-identity}
\end{equation}
The clean Grover subspace of dimension at most two
$\mathcal V=\operatorname{span}\{|\psi\rangle,S_g|\psi\rangle\}$ is
contained in $\mathcal M$. Its transported space
$\mathcal V_B=U\mathcal V$ is invariant under $S_g$ and $Q_B$.
Immediately before every zero reflection, $B^\dagger$ maps the
relevant transported state back to $\mathcal M$. Consequently the old
core reflection and its old clean scratch implementation are valid on
this transported invariant, although terminal odd-degree work is dirty.

For every integer $d\geq1$, the odd-degree good-flag observation and
the even-degree core-zero observation have exactly the same
probabilities for $A$ and $B$. Specifically, if
$p=\langle\psi|\Pi_g|\psi\rangle$ and
$\theta=\arcsin\sqrt p$, both probabilities are
\begin{equation}
 q_d(p)=T_d(\sqrt p)^2
 =\begin{cases}
   \sin^2(d\theta),&d\text{ odd},\\
   \cos^2(d\theta),&d\text{ even}.
  \end{cases}
 \label{eq:conjugation-parity-signal}
\end{equation}
The same equality holds for the complete phase-register distribution
of controlled quantum phase estimation, including an approximate
inverse Fourier circuit applied identically to the phase register.
\end{theorem}
\begin{proof}
Using $B=UA$ and $U^\dagger S_g=S_gU^\dagger$ gives
\begin{equation}
 -B\widetilde S_0B^\dagger S_g
 =-U A\widetilde S_0A^\dagger U^\dagger S_g
 =U(-A\widetilde S_0A^\dagger S_g)U^\dagger.
 \label{eq:conjugation-identity-derivation}
\end{equation}
This is an identity of full operators, not merely an equality of
expectations. Induction on $j$, cancelling every adjacent
$U^\dagger U$, proves the second part of
\eqref{eq:conjugation-global-identity}.

Since $A|e\rangle\in\mathcal M$ and $S_g$ acts only on a core flag,
$\mathcal V\subseteq\mathcal M$. On $\mathcal M$,
\eqref{eq:conjugation-physical-zero-reflection} and invariance under
$A^\dagger$ give
$-A\widetilde S_0A^\dagger=2|\psi\rangle\langle\psi|-I$, so
$\mathcal V$ is invariant under both reflections and their product.
Commutation of $U$ and $S_g$ gives
$\mathcal V_B=\operatorname{span}\{|\phi\rangle,S_g|\phi\rangle\}$
and its invariance. If $x\in\mathcal V$ then
\begin{equation}
 B^\dagger Ux=A^\dagger x\in\mathcal M.
 \label{eq:conjugation-restored-cleanliness}
\end{equation}
Thus work that is dirty during the prepared portion is clean again
after $B^\dagger$, before $S_0$. The existing ladder scratch can be
borrowed from that clean allocation at this point, provided the
emitted inverse clears each borrowed wire and the reflection returns
it to zero. An entangled superposition of phase-control branches does
not change this conclusion: every branch belongs to the same linear
subspace, and the inverse acts unconditionally.

For $d=2k+1$, the two terminal states are
$Q_A^k|\psi\rangle$ and $UQ_A^k|\psi\rangle$. Since
$[U,\Pi_g]=0$, their good-flag probabilities agree. For $d=2k$,
$k\geq1$, the chronological circuit is preparation, $k-1$ Grover
iterations, $S_g$, then the preparation adjoint. Its terminal states
are equal as full vectors:
\begin{equation}
 B^\dagger S_g Q_B^{k-1}B|e\rangle
 =A^\dagger U^\dagger S_gUQ_A^{k-1}A|e\rangle
 =A^\dagger S_gQ_A^{k-1}A|e\rangle\in\mathcal M.
 \label{eq:conjugation-even-state}
\end{equation}
For a vector in $\mathcal M$, core-zero observation is precisely the
rank-one $|e\rangle$ event.

For $0<p<1$, define the orthonormal good and bad components
$|g\rangle=\Pi_g|\psi\rangle/\sqrt p$ and
$|b\rangle=(I-\Pi_g)|\psi\rangle/\sqrt{1-p}$. In that ordered basis,
\begin{equation}
 Q_A\big|_{\mathcal V}
 =\begin{pmatrix}\cos2\theta&\sin2\theta\\
                 -\sin2\theta&\cos2\theta\end{pmatrix},
 \qquad
 Q_A^k|\psi\rangle
 =\sin((2k+1)\theta)|g\rangle+
   \cos((2k+1)\theta)|b\rangle.
 \label{eq:conjugation-grover-matrix}
\end{equation}
The second equality follows by induction using the angle-addition
identities. Odd observation gives $\sin^2((2k+1)\theta)$; even
core-zero amplitude is
\begin{equation}
 \langle\psi|S_g Q_A^{k-1}|\psi\rangle
 =\cos\theta\cos((2k-1)\theta)
  -\sin\theta\sin((2k-1)\theta)
 =\cos(2k\theta).
 \label{eq:conjugation-even-amplitude}
\end{equation}
Since $T_d(\cos x)=\cos(dx)$ and
$\sqrt p=\cos(\pi/2-\theta)$, squaring proves
\eqref{eq:conjugation-parity-signal}. At $p=0$, $Q_A|\psi\rangle
=|\psi\rangle$, and at $p=1$, $Q_A|\psi\rangle=-|\psi\rangle$;
direct evaluation gives respectively odd probabilities $0,1$ and
even probabilities $1,1$. No normalisation of a vanishing component
is used.

For a phase-control qubit $c$, define
$\operatorname{C}_c(V)=|0\rangle\langle0|_c\otimes I+
|1\rangle\langle1|_c\otimes V$. The exact global identity gives
\begin{equation}
 \operatorname{C}_c(Q_B^j)
 =(I_c\otimes U)\operatorname{C}_c(Q_A^j)(I_c\otimes U^\dagger).
 \label{eq:conjugation-controlled}
\end{equation}
This identity is for mathematically controlled full operators. The
actual controlled zero-reflection tape may borrow an extra scratch
wire. It is not assumed globally equal to the controlled version of
the uncontrolled tape on arbitrary dirty work. Let $R_{0,c}$ denote
the actual controlled tape, with its checked restriction
\begin{equation}
 R_{0,c}(|z\rangle_c\otimes x)
 =|z\rangle_c\otimes S_0^z x
 \quad(z\in\{0,1\},\ x\in\mathcal M).
 \label{eq:conjugation-actual-controlled-contract}
\end{equation}
Retain the same actual $R_{0,c}$ in both physical constructions:
\begin{equation}
 K_{A,c}=Z_c(I_c\otimes A)R_{0,c}(I_c\otimes A^\dagger)
                   \operatorname C_c(S_g),\qquad
 K_{B,c}=Z_c(I_c\otimes B)R_{0,c}(I_c\otimes B^\dagger)
                   \operatorname C_c(S_g).
 \label{eq:conjugation-actual-controlled-iterates}
\end{equation}
Commutation of $U$ with $S_g$, and of $I_c\otimes U$ with $Z_c$,
gives the genuine full-tape identity
\begin{equation}
 K_{B,c}=(I_c\otimes U)K_{A,c}(I_c\otimes U^\dagger).
 \label{eq:conjugation-actual-controlled-identity}
\end{equation}
On $\mathcal H_c\otimes\mathcal V$, the intermediate state after
the unconditional $A^\dagger$ has zero work in both control branches,
so \eqref{eq:conjugation-actual-controlled-contract} proves that
$K_{A,c}$ equals $\operatorname C_c(Q_A)$ there. The same reasoning
with \eqref{eq:conjugation-restored-cleanliness} proves equality of
$K_{B,c}$ and $\operatorname C_c(Q_B)$ on
$\mathcal H_c\otimes\mathcal V_B$. These are invariant subspaces,
so repeated controlled iterates implement the required controlled
powers on every reached state, including entangled control branches.

For a register $R$ of controls, each factor in the actual
controlled-power network obeys
\eqref{eq:conjugation-actual-controlled-identity} with $I_R$, so
successive conjugations cancel. Preparation and all actual controlled
powers yield a state equal to $(I_R\otimes U)$ times the $A$ state.
Every exact or approximate Fourier gate acts on $R$ alone and commutes
with $I_R\otimes U$. Therefore this relation holds at terminal phase
observation. For every outcome $y$,
$|y\rangle\langle y|_R\otimes I$ commutes with $I_R\otimes U$,
giving exactly equal outcome probabilities. This also proves equality
of the reduced phase-register states. The minus in
\eqref{eq:conjugation-iterates} is retained throughout: removing it
inside a controlled operation changes relative phases and invalidates
the asserted phase distribution.
\end{proof}

\subsection{Complete resource profiles and the stronger clean baseline}

Let $F=T(U)=T(U^\dagger)$, $P_T=T(P)=T(P^\dagger)$ and let
$h$ be the number of input Hadamards. For additive primitive profiles,
write $\boldsymbol u$, $\boldsymbol u_{\mathrm{inv}}$,
$\boldsymbol p$ and $\boldsymbol h$ for the complete emitted forward
body, inverse body, predicate and input-Hadamard count vectors.
A literal adjoint tape interchanges $T$ and $T^\dagger$, and $S$ and
$S^\dagger$; a separately emitted exact inverse may use an equivalent
decomposition and must retain its own complete profile. The retained
two coherent architectures have equal combined forward/inverse body
$T$ charges. Hence the exact retained and alternative profiles satisfy
\begin{equation}
 \boldsymbol a=\boldsymbol u+\boldsymbol u_{\mathrm{inv}}
                  +\boldsymbol p+\boldsymbol h,
 \qquad \boldsymbol b=\boldsymbol u+\boldsymbol p+\boldsymbol h,
 \qquad a=2F+P_T,\qquad b=F+P_T.
 \label{eq:conjugation-resource-profiles}
\end{equation}
The familiar shorthand $A=2F+P+H$ and $B=F+P+H$ denotes a
composition of emitted profiles; $H$ contributes Clifford operations
and zero $T$ gates. Any claimed improvement in a depth bound must use
the complete removed inverse-body depth profile, not just its $T$ count.
The old larger conservative depth upper bounds may instead be retained
without a depth-improvement claim. The retained
core and conservative peak are unchanged. No further zero-reflection
optimisation, fresh financial precision choice or free arithmetic
register is supplied by this identity.

With $s=T(\widetilde S_0)$, the complete parity shot using $B$ has
\begin{equation}
 t_B(d)=\begin{cases}
  db+(d-1)s/2,&d\text{ odd},\\
  db+(d/2-1)s,&d\text{ even}.
 \end{cases}
 \label{eq:conjugation-parity-resource}
\end{equation}
The same preparation/adjoint count proof as the retained sampler
applies because its old core reflection is justified by
\eqref{eq:conjugation-restored-cleanliness}. Odd terminal work must be
measured and conditionally reset after observation. I retain the
declared full-allocation measurement/reset charge, not a charge for
the core alone. Those operations have zero $T$ cost under the retained
logical architecture but positive measurement, feedback and primitive
operation counts. A coherent downstream consumer requiring clean work
needs a final $U^\dagger$ and must pay its complete profile.

For a fairer clean comparator, cancellation across the preparation
boundary gives the exact full-operator identity
\begin{equation}
 A^\dagger S_g A
 =H^\dagger U^\dagger P^\dagger U S_g U^\dagger P U H
 =H^\dagger U^\dagger P^\dagger S_g P U H.
 \label{eq:conjugation-cross-boundary}
\end{equation}
I denote the clean sampler with every such cancellation performed by
$t_{A,\mathrm{cancel}}(d)$, retaining the same old reflection tape.
There are $(d-1)/2$ disjoint pairs for odd degree and $d/2$ for even
degree, each removing one $U,U^\dagger$ pair. Thus
\begin{equation}
 t_{A,\mathrm{cancel}}(d)=
 \begin{cases}
  t_A(d)-(d-1)F=t_B(d)+F,&d\text{ odd},\\
  t_A(d)-dF=t_B(d),&d\text{ even}.
 \end{cases}
 \label{eq:conjugation-clean-baseline}
\end{equation}
In contrast, $t_A(d)-t_B(d)=dF$ against the old additive template.
Most of that larger difference is therefore a standard cancellation
benefit; it cannot honestly be attributed to a new pricing signal.
These are exact costs of specified constructions, not lower bounds
on every circuit with the same signal.

For a complete nonzero-shot QPE plan, let $M=2^m$ be the phase-grid
size, $R$ the number of independent repetitions, $\tau_{\mathrm{QFT}}$
the complete emitted inverse-Fourier $T$ count, and
$s_c=T(R_{0,c})$ under the old actual controlled-core tape, retaining
its conditional phase and cleanup. Its correct action is required
only on the clean invariant in
\eqref{eq:conjugation-actual-controlled-contract}.
The controlled Grover invocation is emitted as
\begin{equation}
 K_{B,c}
 =Z_c(I_c\otimes B)R_{0,c}
       (I_c\otimes B^\dagger)\operatorname C_c(S_g),
 \qquad T(K_{B,c})=2b+s_c.
 \label{eq:conjugation-controlled-resource}
\end{equation}
Here $\operatorname C(S_g)$ is the charged Clifford controlled-$Z$,
and $Z_c$ supplies the conditional global minus. $B$ and its adjoint
run unconditionally, so there is no assumption that controlling a
financial arithmetic tape is free. On the relevant transported
invariant, control zero gives $BB^\dagger=I$, and control one gives
exactly $Q_B$.
Both branches satisfy the zero-scratch condition before the controlled
reflection on the relevant transported invariant. The complete clean
QPE network, including its initial preparation, is the $B$ network
followed by $U^\dagger$ on the financial register. To see this, the
proof of \eqref{eq:conjugation-actual-controlled-identity} shows that
the $B$ network
has terminal state $(I_R\otimes U)$ times the clean network state;
appending $I_R\otimes U^\dagger$ recovers the latter exactly. The
same circuit results from cross-boundary cancellation in the complete
$A$ network. This is a globally compiled construction. It does not
assert that an isolated invocation of $\operatorname C(Q_A)$, without
its adjacent preparation context, has the reduced cost
$2b+s_c$.

The full retained, alternative and cancellation-aware clean $T$
charges are therefore
\begin{align}
 T_A^{\mathrm{QPE}}
 &=R\{a+(M-1)(2a+s_c)+\tau_{\mathrm{QFT}}\},
 \label{eq:conjugation-retained-qpe-cost}\\
 T_B^{\mathrm{QPE}}
 &=R\{b+(M-1)(2b+s_c)+\tau_{\mathrm{QFT}}\},
 \label{eq:conjugation-dirty-qpe-cost}\\
 T_{A,\mathrm{cancel}}^{\mathrm{QPE}}
 &=R\{a+(M-1)(2b+s_c)+\tau_{\mathrm{QFT}}\}.
 \label{eq:conjugation-clean-qpe-cost}
\end{align}
The differences are
\begin{equation}
 T_A^{\mathrm{QPE}}-T_B^{\mathrm{QPE}}=R(2M-1)F,
 \qquad
 T_{A,\mathrm{cancel}}^{\mathrm{QPE}}-T_B^{\mathrm{QPE}}=RF.
 \label{eq:conjugation-qpe-difference}
\end{equation}
These formulae include one complete preparation, all $M-1$ controlled
Grover invocations and the complete Fourier tape in each repetition.
The full resource record must additionally retain phase Hadamards,
Clifford controlled predicates, the conditional sign, all terminal
observations, full-allocation resets and the old conservative peak
plus phase and Fourier ancillas. Zero-shot prior decisions cost zero
in all three versions. No result above reduces the financial
discretisation allowance, changes the implemented good probability,
improves the QPE confidence theorem or establishes physical runtime
advantage.

\subsection{Controller transfer and the scope of the contribution}

The equality of every $q_d(p)$ in
\eqref{eq:conjugation-parity-signal} is exact within each representation.
For a frozen deterministic controller whose initial information,
confidence allocation and decision policy are unchanged, common
uniform variables therefore couple the $A$ and $B$ observations to
give the same finite history and decisions. A budget-dependent degree
policy must be frozen for this statement: reducing its charged costs
and then rerunning the policy can change which degrees it chooses.
Between different financial representations, separately rounded
probabilities can differ, so the existing equal-signal same-history
corollary remains conditional. The all-degree implementation lemma
may be applied to the new $b_i$ and unchanged $s_i$ to compare their
matched-degree shot charges, but it does not remove that statistical
premise.

The project-specific question is whether the checked financial
representation remains cheaper after every rival receives this
cancellation-aware compilation, identical finite financial
certificates and complete estimator charges. A positive answer would
strengthen the financial implementation evidence by surviving an
established optimisation. It would not constitute a new generic
quantum amplification law or establish theoretical priority for the
conjugation identity.

\begin{thebibliography}{99}
\bibitem[Bayer et al.(2018)]{bayer}
Bayer, C., Siebenmorgen, M. and Tempone, R. (2018).
Smoothing the payoff for efficient computation of basket option prices.
\emph{Quantitative Finance}, 18(3), 491--505.
\url{https://doi.org/10.1080/14697688.2017.1308003}.

\bibitem[Brassard et al.(2002)]{brassard}
Brassard, G., H\o yer, P., Mosca, M. and Tapp, A. (2002).
Quantum amplitude amplification and estimation.
\emph{Contemporary Mathematics}, 305, 53--74.
\url{https://doi.org/10.1090/conm/305/05215}.

\bibitem[Carrera Vazquez and Woerner(2021)]{carreravazquez2021}
Carrera Vazquez, A. and Woerner, S. (2021). Efficient state preparation for quantum amplitude estimation.
\emph{Physical Review Applied}, 15, 034027. \url{https://doi.org/10.1103/PhysRevApplied.15.034027}.
Detailed version consulted: arXiv:2005.07711v1, 15 May 2020.

\bibitem[Curran(1994)]{curran}
Curran, M. (1994). Valuing Asian and portfolio options by conditioning
on the geometric mean price.
\emph{Management Science}, 40(12), 1705--1711.
\url{https://doi.org/10.1287/mnsc.40.12.1705}.

\bibitem[Duchi et al.(2012)]{duchi2012}
Duchi, J. C., Bartlett, P. L. and Wainwright, M. J. (2012). Randomized smoothing for stochastic optimization.
\emph{SIAM Journal on Optimization}, 22(2), 674--701. \url{https://doi.org/10.1137/110831659}.
Detailed version consulted: arXiv:1103.4296v2, 7 April 2012.

\bibitem[Glasserman and Yao(1992)]{glassermanyao1992}
Glasserman, P. and Yao, D. D. (1992). Some guidelines and guarantees for common random numbers.
\emph{Management Science}, 38(6), 884--908. \url{https://doi.org/10.1287/mnsc.38.6.884}.

\bibitem[Marinelli(2024)]{marinelli2024}
Marinelli, C. (2024). On certain representations of pricing functionals.
\emph{Annals of Finance}, 20, 91--127. \url{https://doi.org/10.1007/s10436-024-00438-5}.

\bibitem[Rall and Fuller(2023)]{rallfuller2023} Rall, P. and Fuller, B. (2023). Amplitude estimation from quantum signal processing. \emph{Quantum}, 7, 937. \url{https://doi.org/10.22331/q-2023-03-02-937}.

\bibitem[Rebentrost et al.(2018)]{rebentrost}
Rebentrost, P., Gupt, B. and Bromley, T. R. (2018).
Quantum computational finance: Monte Carlo pricing of financial derivatives.
\emph{Physical Review A}, 98, 022321.
\url{https://doi.org/10.1103/PhysRevA.98.022321}.

\bibitem[Stamatopoulos et al.(2020)]{stamatopoulos2020} Stamatopoulos, N., Egger, D. J., Sun, Y., Zoufal, C., Iten, R., Shen, N. and Woerner, S. (2020). Option pricing using quantum computers. \emph{Quantum}, 4, 291. \url{https://doi.org/10.22331/q-2020-07-06-291}.

\bibitem[Suzuki et al.(2020)]{suzuki2020}
Suzuki, Y., Uno, S., Raymond, R., Tanaka, T., Onodera, T. and Yamamoto, N. (2020).
Amplitude estimation without phase estimation. \emph{Quantum Information Processing}, 19, 75.
\url{https://doi.org/10.1007/s11128-019-2565-2}.
\end{thebibliography}
\end{document}
